\NormalFont\ShortTitle{Luca}

{\MTB Vestea cea bună conform

\par }{\MT Luca


\par }\ChapOne{1}{\PP \VerseOne{1}Deoarece mulți s-au încumetat să facă o relatare despre cele ce s-au împlinit printre noi,

\VS{2}așa cum ni le-au transmis cei care au fost de la început martori oculari și slujitori ai cuvântului,

\VS{3}mi s-a părut bine și mie, după ce am urmărit cu exactitate de la început desfășurarea tuturor lucrurilor, să-ți scriu cu rânduială, preamăritule Teofil,

\VS{4}ca să cunoști certitudinea lucrurilor în care ai fost instruit.


\par }{\PP \VS{5}În zilele lui Irod, împăratul Iudeii, era un preot numit Zaharia, din seminția preoțească a lui Abia. El avea o soție dintre fiicele lui Aaron, și numele ei era Elisabeta.

\VS{6}Amândoi erau neprihăniți înaintea lui Dumnezeu, umblând fără vină în toate poruncile și rânduielile Domnului.

\VS{7}Dar nu au avut niciun copil, pentru că Elisabeta era stearpă, iar ei amândoi erau înaintați în vârstă.


\par }{\PP \VS{8}În timp ce el îndeplinea slujba de preot în fața lui Dumnezeu, în ordinea împărțiriisale,

\VS{9}după obiceiul slujbei de preot, îi revenea sarcina de a intra în templul Domnului și de a aduce tămâie.

\VS{10}Toată mulțimea poporului se ruga afară la ora tămâierii.


\par }{\PP \VS{11}Și i s-a arătat un înger al Domnului, care stătea în picioare la dreapta altarului tămâiei.

\VS{12}Zaharia s-a tulburat când l-a văzut și l-a cuprins frica.

\VS{13}Dar îngerul i-a zis: “Nu te teme, Zaharia, căci cererea ta a fost ascultată. Soția ta, Elisabeta, îți va naște un fiu și îi vei pune numele Ioan.

\VS{14}Vei avea bucurie și veselie, și mulți se vor bucura de nașterea lui.

\VS{15}Căci el va fi mare înaintea Domnului și nu va bea nici vin, nici băutură tare. El va fi plin de Duhul Sfânt, încă din pântecele mamei sale.

\VS{16}El va întoarce pe mulți dintre copiii lui Israel la Domnul Dumnezeul lor.

\VS{17}El va merge înaintea lui în spiritul și cu puterea lui Ilie, “ca să întoarcă inimile părinților la copii”\FTNT{*}{{\FR 1:17 }{\FT{Maleahi 4:6}}} și pe cei neascultători la înțelepciunea celor drepți, ca să pregătească un popor pregătit pentru Domnul.”


\par }{\PP \VS{18}Zaharia a zis îngerului: “Cum pot fi sigur de aceasta? Pentru că eu sunt bătrân, iar soția mea este foarte înaintată în vârstă.”


\par }{\PP \VS{19}Îngerul i-a răspuns: “Eu sunt Gabriel, care stă în fața lui Dumnezeu. Am fost trimis să-ți vorbesc și să-ți aduc această veste bună.

\VS{20}Iată că
\FTNT{†}{{\FR 1:20 }{\FT{“Iată”, de la “ἰδοὺ”, înseamnă a privi, a lua aminte, a observa, a vedea sau a privi. Este adesea folosit ca o interjecție.}}}vei tăcea și nu vei putea vorbi până în ziua în care se vor întâmpla aceste lucruri, pentru că nu ai crezut cuvintele mele, care se vor împlini la timpul lor.”


\par }{\PP \VS{21}Poporul îl aștepta pe Zaharia și se mira că întârzie în templu.

\VS{22}Când a ieșit, nu a putut să le vorbească. Ei și-au dat seama că avusese o viziune în templu. El a continuat să le facă semne și a rămas mut.

\VS{23}Când s-au împlinit zilele de slujbă, a plecat la casa lui.

\VS{24}După aceste zile, Elisabeta, soția sa, a rămas însărcinată și s-a ascuns cinci luni, spunând:

\VS{25}“Așa a făcut Domnul cu mine în zilele în care s-a uitat la mine, ca să-mi îndepărteze ocara printre oameni.”


\par }{\PP \VS{26}În luna a șasea, îngerul Gabriel a fost trimis de Dumnezeu într-o cetate din Galileea, numită Nazaret,

\VS{27}la o fecioară care urma să se mărite cu un bărbat numit Iosif, din casa lui David. Numele fecioarei era Maria.

\VS{28}După ce a intrat, îngerul i-a spus: “Bucură-te, preaiubită! Domnul este cu tine. Binecuvântată ești tu între femei!”


\par }{\PP \VS{29}Dar când l-a văzut, s-a tulburat foarte tare de această vorbă și s-a gândit ce fel de salut ar putea fi acesta.

\VS{30}Îngerul i-a zis: “Nu te teme, Maria, căci ai găsit bunăvoință la Dumnezeu.

\VS{31}Iată că vei concepe în pântecele tău și vei da naștere unui fiu și-i vei pune numele “Isus”.

\VS{32}El va fi mare și va fi numit Fiul Celui Preaînalt. Domnul Dumnezeu îi va da scaunul de domnie al tatălui său, David,

\VS{33}și va domni peste casa lui Iacov pentru totdeauna. Împărăția lui nu va avea sfârșit”.


\par }{\PP \VS{34}Maria a zis îngerului: “Cum se poate, dacă eu sunt fecioară?”


\par }{\PP \VS{35}Îngerul i-a răspuns: “Duhul Sfânt va veni peste tine și puterea Celui Preaînalt te va umbri. De aceea și Sfântul care se va naște din tine se va numi Fiul lui Dumnezeu.

\VS{36}Iată că și Elisabeta, ruda ta, a conceput un fiu la bătrânețe; și aceasta este luna a șasea cu cea care era numită stearpă.

\VS{37}Căci nimic din ceea ce spune Dumnezeu nu este imposibil.”
\FTNT{‡}{{\FR 1:37 }{\FT{2 până la 3 metreți reprezintă aproximativ 20 până la 30 de galoane americane sau 75 până la 115 litri.}}}


\par }{\PP \VS{38}Și Maria a zis: “Iată roaba Domnului; să mi se facă după cuvântul Tău.”

\par }{\PP Apoi îngerul a plecat de la ea.


\par }{\PP \VS{39}În zilele acelea, Maria s-a sculat și s-a dus în grabă în ținutul muntos, într-o cetate din Iuda,

\VS{40}a intrat în casa lui Zaharia și a salutat pe Elisabeta.

\VS{41}Când Elisabeta a auzit salutul Mariei, copilul a sărit în pântecele ei; și Elisabeta s-a umplut de Duhul Sfânt.

\VS{42}Ea a strigat cu glas tare și a zis: “Binecuvântată ești tu între femei și binecuvântat este rodul pântecelui tău!

\VS{43}De ce sunt eu atât de favorizată, încât mama Domnului meu să vină la mine?

\VS{44}Căci iată că, atunci când vocea salutului tău a ajuns la urechile mele, copilul a sărit de bucurie în pântecele meu!

\VS{45}Binecuvântată este cea care a crezut, pentru că se va împlini ceea ce i s-a spus de la Domnul!”


\par }{\PP \VS{46}Maria a spus,

\par }{\Q “Sufletul meu slăvește pe Domnul.


\par }{\QB \VS{47}Duhul meu s-a bucurat în Dumnezeu, Mântuitorul meu,


\par }{\QB \VS{48}Căci a privit la starea umilă a robului Său.

\par }{\Q Căci iată, de acum înainte, toate generațiile mă vor numi binecuvântat.


\par }{\QB \VS{49}Căci Cel ce este puternic a făcut lucruri mari pentru mine.

\par }{\QB Sfânt este numele lui.


\par }{\QB \VS{50}Îndurarea Lui este din generație în generație pentru cei ce se tem de El.


\par }{\Q \VS{51}A dat dovadă de forță cu brațul.

\par }{\QB El i-a împrăștiat pe cei mândri în închipuirea inimilor lor.


\par }{\Q \VS{52}El a dat jos căpeteniile de pe tronurile lor,

\par }{\QB și i-a înălțat pe cei umili.


\par }{\Q \VS{53}El a umplut de bunătăți pe cei flămânzi.

\par }{\QB El i-a trimis pe cei bogați la plimbare cu mâna goală.


\par }{\Q \VS{54}El a ajutat pe Israel, robul Său, ca să-și aducă aminte de îndurare,


\par }{\QB \VS{55}așa cum a vorbit părinților noștri,

\par }{\QB lui Avraam și urmașilor lui în veci.”


\par }{\PP \VS{56}Maria a stat cu ea vreo trei luni și apoi s-a întors la casa ei.


\par }{\PP \VS{57}Și s-a împlinit vremea când Elisabeta trebuia să nască și a născut un fiu.

\VS{58}Vecinii și rudele ei au auzit că Domnul și-a mărit mila față de ea și s-au bucurat împreună cu ea.

\VS{59}În ziua a opta, au venit să taie împrejur copilul; și voiau să-l numească Zaharia, după numele tatălui său.

\VS{60}Mama lui a răspuns: “Nu așa, ci se va numi Ioan”.


\par }{\PP \VS{61}Ei i-au zis: “Nu este nimeni dintre rudele tale care să se numească așa.”

\VS{62}Și i-au făcut semne tatălui său, cum voia să-l cheme.


\par }{\PP \VS{63}A cerut o tăbliță de scris și a scris: “Numele Lui este Ioan”.

\par }{\PP Cu toții s-au mirat.

\VS{64}Gura i s-a deschis îndată și limba i s-a eliberat și a vorbit, binecuvântând pe Dumnezeu.

\VS{65}Frica a cuprins pe toți cei care locuiau în jurul lor și toate aceste cuvinte au fost povestite în tot ținutul muntos al Iudeii.

\VS{66}Toți cei care le auzeau le strângeau în inima lor, spunând: “Ce va fi atunci copilul acesta?” Mâna Domnului era cu el.


\par }{\PP \VS{67}Zaharia, tatăl său, s-a umplut de Duhul Sfânt și a proorocit, zicând,


\par }{\Q \VS{68}“Binecuvântat să fie Domnul, Dumnezeul lui Israel,

\par }{\QB pentru că a vizitat și a răscumpărat pe poporul Său;


\par }{\Q \VS{69}și a ridicat pentru noi un corn de mântuire în casa robului Său David.


\par }{\QB \VS{70}(așa cum a vorbit prin gura sfinților Săi prooroci care au fost din vechime),


\par }{\QB \VS{71}mântuire de vrăjmașii noștri și din mâna tuturor celor care ne urăsc;


\par }{\Q \VS{72}pentru a arăta milă față de părinții noștri,

\par }{\QB să ne aducem aminte de legământul Său sfânt,


\par }{\Q \VS{73}jurământul pe care l-a făcut lui Avraam, tatăl nostru,


\par }{\QB \VS{74}să ne îngăduie ca noi, izbăviți din mâna vrăjmașilor noștri,

\par }{\QB ar trebui să-l slujească fără teamă,


\par }{\QB \VS{75}în sfințenie și dreptate înaintea Lui în toate zilele vieții noastre.


\par }{\Q \VS{76}Și tu, copile, vei fi numit profet al Celui Preaînalt;

\par }{\QB căci tu vei merge înaintea feței Domnului, ca să-i pregătești căile,


\par }{\QB \VS{77}pentru a da cunoștință de mântuire poporului Său prin iertarea păcatelor lor,


\par }{\Q \VS{78}datorită înduratei milostiviri a Dumnezeului nostru,

\par }{\QB prin care ne va vizita zorii de sus,


\par }{\QB \VS{79}să lumineze pe cei ce stau în întuneric și în umbra morții;

\par }{\QB să ne călăuzească picioarele pe calea păcii.”


\par }{\PP \VS{80}Copilul creștea și se întărea în duh și a stat în pustie până în ziua în care a apărut în fața lui Israel.



\par }\Chap{2}{\PP \VerseOne{1}În vremea aceea, Cezar Augustus a dat o poruncă prin care a poruncit ca toată lumea să fie înscrisă.

\VS{2}Aceasta a fost prima înscriere făcută când Quirinius era guvernator al Siriei.

\VS{3}Toți s-au dus să se înscrie, fiecare în cetatea lui.

\VS{4}Și Iosif s-a urcat din Galileea, din cetatea Nazaret, în Iudeea, în cetatea lui David, care se numește Betleem, pentru că era din casa și din familia lui David,

\VS{5}ca să se înscrie împreună cu Maria, care, fiind însărcinată, era făgăduită ca soție, pentru a se căsători cu el.


\par }{\PP \VS{6}Pe când erau ei acolo, a venit ziua în care trebuia să nască.

\VS{7}Ea a născut pe primul ei fiu născut. L-a înfășurat în benzi de pânză și l-a culcat într-o iesle, pentru că nu mai era loc pentru ei în han.


\par }{\PP \VS{8}În același ținut erau niște păstori care stăteau la câmp și vegheau noaptea asupra turmei lor.

\VS{9}Iată că un înger al Domnului stătea lângă ei și slava Domnului strălucea în jurul lor, și ei s-au îngrozit.

\VS{10}Îngerul le-a zis: “Nu vă temeți, căci iată că vă aduc o veste bună de mare bucurie, care va fi pentru tot poporul.

\VS{11}Căci vi s-a născut astăzi, în cetatea lui David, un Mântuitor, care este Hristos Domnul.

\VS{12}Iată semnul care vă este dat: veți găsi un prunc înfășurat în fâșii de pânză, culcat într-o iesle de mâncare.”

\VS{13}Deodată, a fost împreună cu îngerul o mulțime de oștiri cerești care lăudau pe Dumnezeu și spuneau


\par }{\Q \VS{14}“Slavă lui Dumnezeu în locurile prea înalte,

\par }{\QB pe pământ pace, bunăvoință față de oameni.”


\par }{\PP \VS{15}Când îngerii au plecat de la ei în cer, păstorii și-au zis unul altuia: “Să mergem acum la Betleem, ca să vedem ce s-a întâmplat și ce ne-a făcut cunoscut Domnul.”

\VS{16}Au venit în grabă și i-au găsit pe Maria și pe Iosif, iar pruncul era culcat în iesle.

\VS{17}Când l-au văzut, au făcut să se răspândească pe larg cuvântul care le fusese spus despre acest copil.

\VS{18}Toți cei care auzeau se mirau de cele ce le spuseseră păstorii.

\VS{19}Dar Maria a păstrat toate aceste cuvinte, meditându-le în inima ei.

\VS{20}Păstorii s-au întors, slăvind și lăudând pe Dumnezeu pentru toate lucrurile pe care le auziseră și le văzuseră, așa cum le fusese spus.


\par }{\PP \VS{21}După ce s-au împlinit cele opt zile pentru tăierea împrejur a copilului, i s-a pus numele de Isus, care fusese dat de înger înainte de a fi conceput în pântece.


\par }{\PP \VS{22}După ce s-au împlinit zilele de curățire a lor, după legea lui Moise, l-au adus la Ierusalim, ca să-l prezinte Domnului

\VS{23}(după cum este scris în legea Domnului: “Orice bărbat care va deschide pântecele va fi numit sfânt pentru Domnul”),

\VS{24}și să aducă jertfă, după cum este scris în legea Domnului: “O pereche de turturele sau doi pui de porumbel”.


\par }{\PP \VS{25}Iată, era un om în Ierusalim, al cărui nume era Simeon. Omul acesta era neprihănit și evlavios, căutând mângâierea lui Israel, și Duhul Sfânt era peste el.

\VS{26}I se descoperise prin Duhul Sfânt că nu avea să vadă moartea înainte de a-L vedea pe Hristosul Domnului.

\VS{27}El a venit în Duhul Sfânt în templu. Când părinții au adus copilul Isus, ca să facă cu privire la el după obiceiul legii,

\VS{28}el l-a primit în brațe, a binecuvântat pe Dumnezeu și a zis


\par }{\Q \VS{29}“Acum, Stăpâne, dai drumul robului tău,

\par }{\QB după cuvântul Tău, în pace;


\par }{\Q \VS{30}Căci ochii mei au văzut mântuirea Ta,


\par }{\QB \VS{31}pe care ai pregătit-o înaintea ochilor tuturor popoarelor;


\par }{\Q \VS{32}ca o lumină de revelație pentru neamuri,

\par }{\QB și slava poporului Tău Israel.”


\par }{\PP \VS{33}Iosif și mama sa se mirau de cele ce se spuneau despre el.

\VS{34}Simeon i-a binecuvântat și a zis mamei sale, Maria: “Iată, copilul acesta este rânduit pentru căderea și ridicarea multora în Israel și pentru un semn despre care se vorbește.

\VS{35}Da, o sabie va străpunge prin sufletul vostru, ca să fie descoperite gândurile multor inimi.”


\par }{\PP \VS{36}Era o proorociță, Ana, fiica lui Fanuel, din seminția lui Așer (era de o vârstă înaintată, căci trăise cu un bărbat șapte ani de la feciorie,

\VS{37}și era văduvă de vreo optzeci și patru de ani), care nu se depărta de templu, închinându-se zi și noapte cu post și cu cereri.

\VS{38}Subind chiar în ceasul acela, ea aducea mulțumiri Domnului și vorbea despre el tuturor celor care așteptau răscumpărarea în Ierusalim.


\par }{\PP \VS{39}După ce au împlinit tot ce era după Legea Domnului, s-au întors în Galileea, în cetatea lor, Nazaret.

\VS{40}Copilul creștea și se întărea în spirit, fiind plin de înțelepciune, și harul lui Dumnezeu era peste el.


\par }{\PP \VS{41}Părinții lui se duceau în fiecare an la Ierusalim, la sărbătoarea Paștelui.

\VS{42}Când a împlinit vârsta de doisprezece ani, s-au suit la Ierusalim, după obiceiul sărbătorii;

\VS{43}și, după ce au împlinit zilele, pe când se întorceau, băiatul Isus a rămas în Ierusalim. Iosif și mama lui nu știau acest lucru,

\VS{44}dar, presupunând că este în companie, au plecat o zi de drum; și l-au căutat printre rudele și cunoscuții lor.

\VS{45}Cum nu l-au găsit, s-au întors la Ierusalim, în căutarea lui.

\VS{46}După trei zile, l-au găsit în templu, stând în mijlocul învățătorilor, atât pentru a-i asculta, cât și pentru a le pune întrebări.

\VS{47}Toți cei care îl ascultau erau uimiți de înțelegerea și de răspunsurile lui.

\VS{48}Când l-au văzut, au rămas uimiți, iar mama lui i-a zis: “Fiule, de ce te-ai purtat așa cu noi? Iată că tatăl tău și cu mine te așteptam cu nerăbdare”.


\par }{\PP \VS{49}El le-a zis: “{\WJ{Pentru ce Mă căutați? Nu știați că trebuie să fiu în casa Tatălui Meu?”
}}

\VS{50}Ei nu au înțeles cuvântul pe care li-l spunea.

\VS{51}Și, coborând cu ei, a venit la Nazaret. El era supus lor, iar mama lui păstra în inima ei toate aceste cuvinte.

\VS{52}Și Isus creștea în înțelepciune și în statură și în favoare înaintea lui Dumnezeu și a oamenilor.



\par }\Chap{3}{\PP \VerseOne{1}În al cincisprezecelea an al domniei lui Tiberiu Cezar, când Ponțiu Pilat era guvernator al Iudeii, Irod era tetrarh al Galileii, iar fratele său Filip era tetrarh al ținutului Ituraei și al Traconitei și Lisanias tetrarh al Abilenei,

\VS{2}în timpul marii preoții a lui Ana și Caiafa, cuvântul lui Dumnezeu a venit lui Ioan, fiul lui Zaharia, în pustie.

\VS{3}El a venit în toată regiunea din jurul Iordanului, predicând botezul pocăinței pentru iertarea păcatelor.

\VS{4}După cum este scris în cartea cuvintelor profetului Isaia,

\par }{\Q “Glasul cuiva care strigă în pustiu,

\par }{\QB 'Pregătiți calea Domnului.

\par }{\Q Fă-i cărările drepte.


\par }{\QB \VS{5}Orice vale va fi umplută.

\par }{\Q Fiecare munte și deal va fi doborât.

\par }{\QB Cele strâmbe vor deveni drepte,

\par }{\QB și căile aspre să fie netede.


\par }{\Q \VS{6}Toată făptura va vedea mântuirea lui Dumnezeu.”


\par }{\PP \VS{7}De aceea a zis mulțimilor care ieșeau să fie botezate de El: “Neam de vipere, cine v-a pus în gardă să fugiți de mânia viitoare?

\VS{8}De aceea, dați roade demne de pocăință și nu începeți să spuneți între voi: “Avem ca tată pe Avraam!” Căci vă spun că Dumnezeu poate să-i ridice lui Avraam copii din aceste pietre!”.

\VS{9}Chiar și acum toporul se află la rădăcina copacilor. Așadar, orice copac care nu produce fructe bune este tăiat și aruncat în foc.”


\par }{\PP \VS{10}Și mulțimile L-au întrebat: “Ce trebuie să facem?”


\par }{\PP \VS{11}El le-a răspuns: “Cine are două haine, să le dea celui care nu are. Cine are de mâncare, să facă la fel.”


\par }{\PP \VS{12}Au venit și vameșii să fie botezați și I-au zis: “Învățătorule, ce trebuie să facem?”


\par }{\PP \VS{13}El le-a zis: “Nu culegeți mai mult decât ceea ce vi se cuvine”.


\par }{\PP \VS{14}Și soldații l-au întrebat și ei: “Și noi? Ce trebuie să facem noi?”

\par }{\PP El le-a spus: “Nu luați de la nimeni prin violență și nu acuzați pe nimeni pe nedrept. Fiți mulțumiți cu salariile voastre”.


\par }{\PP \VS{15}Pe când norodul era în așteptare și toți oamenii se gândeau în inimile lor dacă nu cumva Ioan este Hristosul,

\VS{16}Ioan le-a răspuns tuturor: “Eu vă botez cu apă, dar vine Cel mai puternic decât mine, căruia nu sunt vrednic să-i dezleg cureaua sandalelor. El vă va boteza în Duhul Sfânt și în foc.

\VS{17}În mâna Lui este vânturatorul, va curăța bine aria și va strânge grâul în hambarul Său, dar va arde neghina cu foc nestins.”


\par }{\PP \VS{18}Apoi, cu multe alte îndemnuri, propovăduia poporului vestea cea bună.

\VS{19}Dar Irod tetrarhul, fiind mustrat de el pentru Irodiada, soția fratelui său, și pentru toate relele pe care le făcuse Irod,

\VS{20}a adăugat la toate acestea și aceasta: a închis pe Ioan în temniță.


\par }{\PP \VS{21}Și când tot poporul era botezat, Isus, botezat și el, se ruga. Cerul s-a deschis

\VS{22}și Duhul Sfânt s-a pogorât peste el în chip trupesc, ca un porumbel, iar din cer s-a auzit un glas care zicea: “Tu ești Fiul Meu preaiubit. În tine sunt binevoitor”.


\par }{\PP \VS{23}Isus însuși, când a început să învețe, era în vârstă de vreo treizeci de ani, fiind fiul lui Iosif, fiul lui Heli,

\VS{24}fiul lui Matei, fiul lui Levi, fiul lui Melchi, fiul lui Jannai, fiul lui Iosif,

\VS{25}fiul lui Matatia, fiul lui Amos, fiul lui Naum, fiul lui Esli, fiul lui Naggai,

\VS{26}fiul lui Maat, fiul lui Matatia, fiul lui Semein, fiul lui Iosif, fiul lui Iuda,

\VS{27}fiul lui Joanan, fiul lui Reșea, fiul lui Zorobabel, fiul lui Șealtiel, fiul lui Neri,

\VS{28}fiul lui Melchi, fiul lui Addi, fiul lui Cosam, fiul lui Elmodam, fiul lui Er,

\VS{29}fiul lui Iosif, fiul lui Eliezer, fiul lui Jorim, fiul lui Matei, fiul lui Levi,

\VS{30}fiul lui Simeon, fiul lui Iuda, fiul lui Iosif, fiul lui Iona, fiul lui Eliachim,

\VS{31}fiul lui Melea, fiul lui Menan, fiul lui Matata, fiul lui Natan, fiul lui David,

\VS{32}fiul lui Isai, fiul lui Obed, fiul lui Boaz, fiul lui Salmon, fiul lui Nahșon,

\VS{33}fiul lui Aminadab, fiul lui Aram, fiul lui Hezron, fiul lui Pereț, fiul lui Iuda,

\VS{34}fiul lui Iacov, fiul lui Isaac, fiul lui Avraam, fiul lui Terah, fiul lui Nahor,

\VS{35}fiul lui Serug, fiul lui Reu, fiul lui Peleg, fiul lui Eber, fiul lui Șelah,

\VS{36}fiul lui Cainan, fiul lui Arfaxad, fiul lui Sem, fiul lui Noe, fiul lui Lameh,

\VS{37}fiul lui Matusalem, fiul lui Enoh, fiul lui Iared, fiul lui Mahalaleel, fiul lui Cainan,

\VS{38}fiul lui Enos, fiul lui Set, fiul lui Adam, fiul lui Dumnezeu.



\par }\Chap{4}{\PP \VerseOne{1}Isus, plin de Duhul Sfânt, s-a întors de la Iordan și a fost dus de Duhul în pustiu

\VS{2}timp de patruzeci de zile, fiind ispitit de diavol. În acele zile nu a mâncat nimic. După aceea, când acestea s-au încheiat, i-a fost foame.


\par }{\PP \VS{3}Diavolul I-a zis: “Dacă ești Fiul lui Dumnezeu, poruncește ca piatra aceasta să se facă pâine.”


\par }{\PP \VS{4}Isus i-a răspuns: “{\WJ{Este scris: “Nu numai cu pâine va trăi omul, ci cu orice cuvânt al lui Dumnezeu.””
}}


\par }{\PP \VS{5}Diavolul, ducându-l pe un munte înalt, i-a arătat într-o clipă toate împărățiile lumii.

\VS{6}Diavolul i-a zis: “Îți voi da toate aceste autorități și slava lor, căci mi-au fost date mie și le dau oricui vreau.

\VS{7}Așadar, dacă te vei închina înaintea mea, toate acestea vor fi ale tale.”


\par }{\PP \VS{8}Isus i-a răspuns: “{\WJ{Du-te de aici, Satano! Căci este scris: “Să te închini Domnului, Dumnezeului tău, și numai lui să-i slujești!”.
}}


\par }{\PP \VS{9}L-a dus la Ierusalim, l-a așezat pe vârful templului și i-a zis: “Dacă ești Fiul lui Dumnezeu, aruncă-te jos de aici,

\VS{10}căci este scris,

\par }{\Q “El va pune pe îngerii Săi să vă păzească, ca să vă păzească.


\par }{\PP \VS{11}și,

\par }{\Q 'Pe mâinile lor te vor purta,

\par }{\QB ca nu cumva să-ți lovești piciorul de o piatră.””


\par }{\PP \VS{12}Isus, răspunzând, i-a zis: “{\WJ{S-a spus: “Să nu ispitești pe Domnul Dumnezeul tău”.”
}}


\par }{\PP \VS{13}După ce a ispitit pe toți, diavolul a plecat de la el până la o altă dată.


\par }{\PP \VS{14}Isus s-a întors în Galileea cu puterea Duhului Sfânt, și vestea despre El s-a răspândit în tot ținutul din jur.

\VS{15}El învăța în sinagogile lor, fiind slăvit de toți.


\par }{\PP \VS{16}Și a venit la Nazaret, unde fusese crescut. A intrat, după obiceiul său, în sinagogă, în ziua de Sabat, și s-a ridicat în picioare să citească.

\VS{17}I s-a înmânat cartea profetului Isaia. A deschis cartea și a găsit locul în care era scris,


\par }{\Q \VS{18}{\WJ{“Duhul Domnului este peste mine,
}}

\par }{\QB {\WJ{pentru că El m-a uns să vestesc săracilor vestea cea bună.
}}

\par }{\Q {\WJ{El m-a trimis să-i vindec pe cei cu inimazdrobită,
}}

\par }{\QB {\WJ{pentru a vesti eliberarea celor prizonieri,
}}

\par }{\QB {\WJ{redarea vederii orbilor,
}}

\par }{\QB {\WJ{pentru a-i elibera pe cei care sunt zdrobiți,
}}


\par }{\QB \VS{19}{\WJ{și să vestească anul de grație al Domnului.”
}}


\par }{\PP \VS{20}A închis cartea, a dat-o înapoi slujbașului și s-a așezat. Ochii tuturor celor din sinagogă erau ațintiți asupra lui.

\VS{21}El a început să le spună: “{\WJ{Astăzi s-a împlinit în auzul vostru Scriptura aceasta.”
}}


\par }{\PP \VS{22}Toți mărturiseau despre el și se mirau de cuvintele pline de har care ieșeau din gura lui și ziceau: “Nu este acesta fiul lui Iosif?”


\par }{\PP \VS{23}El le-a zis: “{\WJ{Fără îndoială că îmi veți spune proverbul acesta: “Doctore, vindecă-te! Tot ce am auzit că s-a făcut în Capernaum, fă și aici, în orașul tău natal!””.
}}

\VS{24}El a zis: “{\WJ{Adevărat vă spun că niciun profet nu este acceptat în orașul său natal.
}}

\VS{25}{\WJ{Dar, adevărat vă spun că erau multe văduve în Israel pe vremea lui Ilie, când cerul a fost închis trei ani și șase luni, când a venit o foamete mare în toată țara.
}}

\VS{26}{\WJ{Ilie nu a fost trimis la niciuna dintre ele, decât la Sarepta, în ținutul Sidonului, la o femeie văduvă.
}}

\VS{27}{\WJ{Pe vremea profetului Elisei, în Israel, erau mulți leproși, dar niciunul dintre ei nu a fost curățat, cu excepția lui Naaman, sirianul.”
}}


\par }{\PP \VS{28}Toți s-au umplut de mânie în sinagogă, când auzeau aceste lucruri.

\VS{29}S-au sculat, l-au aruncat afară din cetate și l-au dus pe vârful dealului pe care era zidită cetatea lor, ca să-l arunce de pe stâncă.

\VS{30}Dar el, trecând prin mijlocul lor, și-a văzut de drum.


\par }{\PP \VS{31}S-a coborât la Capernaum, o cetate din Galileea. Îi învăța în ziua de Sabat;

\VS{32}și ei erau uimiți de învățătura Lui, căci cuvântul Lui era cu autoritate.

\VS{33}În sinagogă era un om care avea un duh de demon necurat; și a strigat cu glas tare,

\VS{34}zicând: “Ah! ce avem noi de-a face cu tine, Isus din Nazaret? Ai venit să ne distrugi? Eu știu cine ești: Sfântul lui Dumnezeu!”.


\par }{\PP \VS{35}Isus l-a mustrat, zicând: “{\WJ{Taci și ieși de la el!” După ce l-a aruncat jos în mijlocul lor, demonul a ieșit din el, fără să-i facă vreun rău.
}}


\par }{\PP \VS{36}Toți au fost cuprinși de uimire și vorbeau între ei, zicând: “Ce este cuvântul acesta? Pentru că el poruncește cu autoritate și putere duhurilor necurate și ele ies!”

\VS{37}Vestea despre el s-a răspândit în toate locurile din regiunea înconjurătoare.


\par }{\PP \VS{38}S-a sculat din sinagogă și a intrat în casa lui Simon. Soacra lui Simon era bolnavă de o febră mare și l-au rugat să o ajute.

\VS{39}El a stat lângă ea și a mustrat febra, iar aceasta a lăsat-o. Imediat s-a sculat și i-a servit.

\VS{40}Pe când apunea soarele, toți cei care aveau pe cineva bolnav de diferite boli i-au adus la el; el și-a pus mâinile peste fiecare dintre ei și i-a vindecat.

\VS{41}Din mulți ieșeau și demoni, care strigau și ziceau: “Tu ești Hristosul, Fiul lui Dumnezeu!” Reproșându-le, nu le-a permis să vorbească, pentru că știau că el era Hristosul.


\par }{\PP \VS{42}Când s-a făcut ziuă, a plecat și s-a dus într-un loc pustiu; și mulțimile Îl căutau, veneau la El și se țineau de El, ca să nu plece de lângă ele.

\VS{43}Dar el le-a spus: “{\WJ{Trebuie să propovăduiesc vestea bună a Împărăției lui Dumnezeu și în celelalte cetăți. De aceea am fost trimis”.
}}

\VS{44}El predica în sinagogile din Galileea.



\par }\Chap{5}{\PP \VerseOne{1}Și, pe când mulțimea îl strângea și asculta cuvântul lui Dumnezeu, el stătea lângă lacul Ghenareț.

\VS{2}A văzut două corăbii care stăteau lângă lac, dar pescarii ieșiseră din ele și își spălau plasele.

\VS{3}A intrat într-una dintre bărci, care era a lui Simon, și l-a rugat să se depărteze puțin de țărm. S-a așezat și a învățat mulțimile din barcă.


\par }{\PP \VS{4}După ce a isprăvit de vorbit, a zis lui Simon: “Scoateți
{\WJ{în larg și aruncați mrejele ca să prindeți.”
}}


\par }{\PP \VS{5}Simon I-a răspuns: “Învățătorule, am lucrat toată noaptea și n-am prins nimic; dar, la cuvântul Tău, voi da drumul la mreajă.”

\VS{6}După ce au făcut acest lucru, au prins o mare mulțime de pești, iar plasa lor se rupea.

\VS{7}Ei au făcut semn partenerilor lor din cealaltă barcă să vină să-i ajute. Aceștia au venit și au umplut ambele bărci, astfel încât acestea au început să se scufunde.

\VS{8}Dar Simon Petru, când a văzut, a căzut la genunchii lui Isus și a zis: “Doamne, depărtează-te de la mine, căci sunt un om păcătos!”.

\VS{9}Căci era uimit, el și toți cei care erau cu el, de capturile de pește pe care le prinseseră;

\VS{10}la fel și Iacov și Ioan, fiii lui Zebedei, care erau asociați cu Simon.

\par }{\PP Isus i-a spus lui Simon:
{\WJ{“Nu te teme. De acum înainte vei prinde oameni vii”.
}}


\par }{\PP \VS{11}După ce și-au adus corăbiile la țărm, au lăsat totul și au mers după El.


\par }{\PP \VS{12}Pe când era într-una din cetăți, iată că era un om plin de lepră. Când L-a văzut pe Isus, a căzut cu fața la pământ și L-a implorat, zicând: “Doamne, dacă vrei, poți să mă cureți.”


\par }{\PP \VS{13}El a întins mâna și l-a atins, zicând:
{\WJ{“Vreau. Să fiu curățat”.
}}

\par }{\PP Imediat lepra l-a părăsit.

\VS{14}I-a poruncit să nu spună nimănui:
{\WJ{“Ci du-te, arată-te preotului și adu jertfă pentru curățirea ta, după cum a poruncit Moise, ca mărturie pentru ei.”
}}


\par }{\PP \VS{15}Dar vestea despre El s-a răspândit mult mai mult și mulțimi mari se adunau să-L asculte și să fie vindecate de bolile lor.

\VS{16}Dar el s-a retras în pustie și se ruga.


\par }{\PP \VS{17}Într-una din zilele acelea, El învăța pe oameni; și stăteau acolo farisei și învățători ai Legii, veniți din toate satele Galileii, din Iudeea și din Ierusalim. Puterea Domnului era cu el ca să-i vindece.

\VS{18}Iată că niște oameni au adus un paralitic pe un pat de campanie și căutau să-l aducă înăuntru ca să-l pună înaintea lui Isus.

\VS{19}Nu au găsit cum să-l aducă înăuntru din cauza mulțimii, s-au urcat pe acoperișul casei și l-au coborât prin țigle cu patul lui în mijlocul ei, în fața lui Isus.

\VS{20}Văzând credința lor, Isus i-a spus:
{\WJ{“Omule, păcatele tale îți sunt iertate.”
}}


\par }{\PP \VS{21}Cărturarii și fariseii au început să se certe, zicând: “Cine este acesta care vorbește de hulă? Cine poate ierta păcatele, dacă nu numai Dumnezeu?”


\par }{\PP \VS{22}Dar Isus, înțelegând gândurile lor, le-a răspuns: “{\WJ{De ce gândiți așa în inimile voastre?
}}

\VS{23}{\WJ{Ce este mai ușor să spui: “Păcatele tale îți sunt iertate” sau să spui: “Scoală-te și umblă”?
}}

\VS{24}{\WJ{Dar ca să știți că Fiul Omului are putere pe pământ să ierte păcatele”, i-a
}}spus paraliticului: “{\WJ{Îți spun: scoală-te, ia-ți patul și du-te la casa ta”.
}}


\par }{\PP \VS{25}Și îndată s-a sculat în fața lor, a luat ce avea pe el și a plecat la casa lui, slăvind pe Dumnezeu.

\VS{26}Toată lumea a fost cuprinsă de uimire și a slăvit pe Dumnezeu. S-au umplut de teamă, spunând: “Am văzut astăzi lucruri ciudate.”


\par }{\PP \VS{27}După aceea, ieșind afară, a văzut un vameș, numit Levi, care ședea la fisc, și i-a zis:
{\WJ{“Vino după mine!”
}}


\par }{\PP \VS{28}A lăsat totul, s-a sculat și l-a urmat.

\VS{29}Levi i-a făcut un mare ospăț în casa lui. Era o mare mulțime de vameși și de alți oameni care se odihneau cu ei.

\VS{30}Cărturarii și fariseii cârteau împotriva ucenicilor Lui, zicând: “De ce mâncați și beți cu vameșii și cu păcătoșii?”


\par }{\PP \VS{31}Isus le-a răspuns: “{\WJ{Cei sănătoși nu au nevoie de doctor, dar cei bolnavi au nevoie de doctor.
}}

\VS{32}{\WJ{Eu n-am venit să chem pe cei drepți, ci pe păcătoși la pocăință.”
}}


\par }{\PP \VS{33}Ei I-au zis: “Pentru ce ucenicii lui Ioan postesc și se roagă des, și ucenicii fariseilor, iar ai voștri mănâncă și beau?”


\par }{\PP \VS{34}El le-a zis: “{\WJ{Puteți voi să faceți pe prietenii mirelui să postească, când mirele este cu ei?
}}

\VS{35}{\WJ{Dar vor veni zile când mirele va fi luat de la ei. Atunci vor posti în acele zile.”
}}


\par }{\PP \VS{36}Și le-a spus și o pildă.
{\WJ{“Nimeni nu pune o bucată dintr-o haină nouă pe o haină veche, căci altfel va rupe haina nouă și, de asemenea, bucata din haina nouă nu se va potrivi cu cea veche.
}}

\VS{37}{\WJ{Nimeni nu pune vin nou în piei de vin vechi, căci altfel vinul nou va sparge piei, se va vărsa și piei vor fi distruse.
}}

\VS{38}{\WJ{Dar vinul nou trebuie pus în piei de vin proaspăt, și amândouă se păstrează.
}}

\VS{39}Nimeni, după ce a
{\WJ{băut vin vechi, nudorește imediat vin nou, pentru că spune: “Cel vechi este mai bun”.”
}}



\par }\Chap{6}{\PP \VerseOne{1}În al doilea Sabat după cel dintâi, el mergea prin lanurile de grâu. Discipolii lui culegeau căpățânile de grâu și mâncau, frecându-le în mâini.

\VS{2}Dar unii dintre farisei le-au spus: “De ce faceți ceea ce nu este permis să faceți în ziua de Sabat?”


\par }{\PP \VS{3}Isus le-a răspuns: “{\WJ{N-ați citit ce a făcut David, când a flămânzit, el și cei ce erau cu el,
}}

\VS{4}{\WJ{cum a intrat în casa lui Dumnezeu, a luat și a mâncat pâinea de sărbătoare și a dat și celor ce erau cu el, care nu se poate mânca decât numai pentru preoți?”
}}

\VS{5}El le-a răspuns:
{\WJ{“Fiul Omului este stăpânul Sabatului”.
}}


\par }{\PP \VS{6}În altă zi de Sabat, a intrat în sinagogă și învăța. Era acolo un om, căruia i se uscase mâna dreaptă.

\VS{7}Cărturarii și fariseii îl urmăreau, ca să vadă dacă nu cumva vindecă în Sabat, ca să găsească o acuzație împotriva lui.

\VS{8}Dar El le cunoștea gândurile; și a zis omului care avea mâna uscată:
{\WJ{“Ridică-te și stai în mijloc”. El s-a ridicat și a stat în picioare.
}}

\VS{9}Atunci Isus le-a zis:
{\WJ{“Vă voi întreba ceva: “Este îngăduit în Sabat să faci bine sau să faci rău? Să salvezi o viață, sau să ucizi?”
}}

\VS{10}S-a uitat la toți și a zis omului:
{\WJ{“Întinde-ți mâna”. El a făcut-o, și mâna i-a fost redată la fel de sănătoasă ca și cealaltă.
}}

\VS{11}Dar ei erau plini de furie și vorbeau între ei despre ce ar putea să-i facă lui Isus.


\par }{\PP \VS{12}În zilele acelea, s-a dus pe munte să se roage și a stat toată noaptea în rugăciune către Dumnezeu.

\VS{13}Când s-a făcut ziuă, și-a chemat discipolii și a ales dintre ei doisprezece, pe care i-a numit și apostoli:

\VS{14}Simon, pe care l-a numit și Petru; Andrei, fratele său; Iacov; Ioan; Filip; Bartolomeu;

\VS{15}Matei; Toma; Iacov, fiul lui Alfeu; Simon, numit Zelota;

\VS{16}Iuda, fiul lui Iacov; și Iuda Iscarioteanul, care a devenit și el trădător.


\par }{\PP \VS{17}S-a coborât cu ei și a stat pe un loc drept, cu o mulțime de ucenici ai Lui și cu un mare număr de oameni din toată Iudeea, din Ierusalim și de pe țărmul mării Tirului și al Sidonului, care veniseră să-L asculte și să se vindece de bolile lor,

\VS{18}precum și de cei care erau tulburați de duhuri necurate; și se vindecau.

\VS{19}Toată mulțimea căuta să se atingă de El, căci ieșea din El o putere și îi vindeca pe toți.


\par }{\PP \VS{20}Și-a ridicat ochii spre ucenicii Săi și a zis:

\par }{\Q {\WJ{“Ferice de voi, cei săraci!
}}

\par }{\QB {\WJ{pentru că Împărăția lui Dumnezeu este a ta.
}}


\par }{\Q \VS{21}{\WJ{Ferice de voi, care sunteți flămânzi acum!
}}

\par }{\QB {\WJ{pentru că vei fi umplut.
}}

\par }{\Q {\WJ{Ferice de voi care plângeți acum,
}}

\par }{\QB {\WJ{pentru că vei râde.
}}


\par }{\Q \VS{22}{\WJ{Ferice de voi, când vă vor urî oamenii, când vă vor respinge, când vă vor batjocori și când vor arunca numele vostru ca fiind rău, din pricina Fiului Omului!
}}


\par }{\QB \VS{23}{\WJ{Bucurați-vă în ziua aceea și săriți de bucurie, căci iată că răsplata voastră este mare în ceruri, pentru că părinții lor au făcut același lucru cu profeții.
}}


\par }{\BB \par }
{\Q \VS{24}“{\WJ{Dar vai de voi, care sunteți bogați!
}}

\par }{\QB {\WJ{Căci ați primit mângâierea voastră.
}}


\par }{\Q \VS{25}{\WJ{Vai de voi, cei ce sunteți plini acum!
}}

\par }{\QB {\WJ{pentru că vă va fi foame.
}}

\par }{\Q {\WJ{Vai de voi care râdeți acum,
}}

\par }{\QB {\WJ{pentru că veți jeli și veți plânge.
}}


\par }{\Q \VS{26}{\WJ{Vai, când oamenii vorbesc de bine despre tine!
}}

\par }{\QB {\WJ{căci părinții lor au făcut același lucru cu profeții mincinoși.
}}


\par }{\PP \VS{27}{\WJ{Dar Eu vă spun vouă, celor ce auziți: iubiți pe vrăjmașii voștri, faceți bine celor ce vă urăsc,
}}

\VS{28}{\WJ{binecuvântați pe cei ce vă blestemă și rugați-vă pentru cei ce vă maltratează.
}}

\VS{29}Celui ce
{\WJ{vă lovește peste obraz, oferiți-i și celălalt obraz; și celui ce vă ia haina, nu-i reține și haina.
}}

\VS{30}{\WJ{Dați oricui vă cere și nu cereți celui care vă ia bunurile să vi le dea înapoi.
}}


\par }{\PP \VS{31}“Fă-le și tu
{\WJ{cum vrei să facă oamenii cu tine.
}}


\par }{\PP \VS{32}{\WJ{Dacă iubești pe cei ce te iubesc, ce merite ai tu? Căci chiar și păcătoșii îi iubesc pe cei care îi iubesc.
}}

\VS{33}{\WJ{Dacă faceți bine celor ce vă fac bine, ce merite aveți voi? Căci și păcătoșii fac la fel.
}}

\VS{34}{\WJ{Dacă dați cu împrumut celor de la care sperați să primiți, ce merite aveți? Chiar și păcătoșii împrumută păcătoșilor, pentru a primi înapoi la fel de mult.
}}

\VS{35}{\WJ{Dar iubiți pe vrăjmașii voștri, faceți binele și împrumutați, fără să așteptați nimic în schimb; și răsplata voastră va fi mare și veți fi copiii Celui Preaînalt, căci El este binevoitor față de cei nerecunoscători și răi.
}}


\par }{\Q \VS{36}{\WJ{“De aceea fiți milostivi,
}}

\par }{\QB {\WJ{așa cum și Tatăl vostru este milostiv.
}}


\par }{\Q \VS{37}{\WJ{Nu judecați,
}}

\par }{\QB {\WJ{și nu vei fi judecat.
}}

\par }{\Q {\WJ{Nu condamnați,
}}

\par }{\QB {\WJ{și nu vei fi condamnat.
}}

\par }{\Q {\WJ{Eliberați,
}}

\par }{\QB {\WJ{și vei fi eliberat.
}}


\par }{\PP \VS{38}{\WJ{“Dați și vi se va da; o măsură bună, presată, zdruncinată și plină, vi se va da. Căci cu aceeași măsură cu care măsurați, vi se va măsura înapoi.”
}}


\par }{\PP \VS{39}Și le-a spus o pildă.
{\WJ{“Poate orbul să călăuzească pe orb? Nu vor cădea amândoi într-o groapă?
}}

\VS{40}{\WJ{Un ucenic nu este mai presus de învățătorul său, dar fiecare, când va fi pe deplin instruit, va fi ca învățătorul său.
}}

\VS{41}{\WJ{De ce vedeți paiul de paie care este în ochiul fratelui vostru, dar nu vă gândiți la bârna care este în ochiul vostru?
}}

\VS{42}{\WJ{Sau cum poți să-i spui fratelui tău: “Frate, lasă-mă să scot paiul din ochiul tău”, când tu însuți nu vezi bârna care este în ochiul tău? Ipocritule! Îndepărtează mai întâi bârna din ochiul tău și apoi vei putea vedea limpede ca să îndepărtezi paiul care este în ochiul fratelui tău.
}}


\par }{\PP \VS{43}{\WJ{“Căci nu este pom bun care să dea roade putrede, nici pom putred care să dea roade bune.
}}

\VS{44}{\WJ{Căci fiecare pom se cunoaște după rodul său. Căci oamenii nu culeg smochine din spini, nici struguri din tufișuri de mărăcini.
}}

\VS{45}{\WJ{Omul bun scoate din comoara bună a inimii sale ceea ce este bun, iar omul rău scoate din comoara rea a inimii sale ceea ce este rău, căci din belșugul inimii vorbește gura lui.
}}


\par }{\PP \VS{46}{\WJ{“Pentru ce Mă strigați: “Doamne, Doamne!” și nu faceți ce vă spun Eu?
}}

\VS{47}{\WJ{Oricine vine la Mine, care ascultă cuvintele Mele și le face, vă voi arăta cum este el.
}}

\VS{48}{\WJ{El este ca un om care construiește o casă, care a săpat, a mers adânc și a pus temelia pe stâncă. Când s-a ridicat un potop, pârâul s-a izbit de casa aceea și nu a putut să o zdruncine, pentru că era întemeiată pe stâncă.
}}

\VS{49}{\WJ{Darcel care aude și nu face, se aseamănă cu un om care a zidit o casă pe pământ fără temelie, împotriva căreia a izbucnit pârâul și îndată a căzut; și ruina acelei case a fost mare.”
}}



\par }\Chap{7}{\PP \VerseOne{1}După ce a terminat de vorbit în auzul poporului, a intrat în Capernaum.

\VS{2}Un slujitor al unui centurion, care îi era foarte drag, era bolnav și era la un pas de moarte.

\VS{3}Când a auzit despre Isus, a trimis la el pe bătrânii iudeilor, rugându-l să vină și să-i salveze slujitorul.

\VS{4}Când au ajuns la Isus, l-au rugat insistent, spunând: “Este vrednic ca tu să faci asta pentru el,

\VS{5}pentru că el iubește neamul nostru și a zidit sinagoga noastră pentru noi”.

\VS{6}Isus a mers cu ei. Când nu mai era acum departe de casă, centurionul a trimis la el niște prieteni, spunându-i: “Doamne, nu te îngrijora, căci nu sunt vrednic să intri sub acoperișul meu.

\VS{7}De aceea nici eu nu m-am considerat vrednic să vin la tine; dar spune cuvântul și robul meu va fi vindecat.

\VS{8}Căci și eu sunt un om pus sub autoritate, având sub mine soldați. Îi spun acestuia: “Du-te!” și el se duce; altuia: “Vino!” și vine; iar robului meu: “Fă aceasta” și el o face”.”


\par }{\PP \VS{9}Isus, auzind acestea, s-a mirat de el și, întorcându-se, a zis mulțimii care-L urma: “{\WJ{Vă spun că n-am găsit o credință atât de mare, nici în Israel.”
}}

\VS{10}Cei care fuseseră trimiși, întorcându-se la casă, au constatat că servitorul care fusese bolnav era sănătos.


\par }{\PP \VS{11}Curând după aceea, s-a dus într-o cetate numită Nain. Mulți dintre discipolii Săi, împreună cu o mare mulțime, au mers cu El.

\VS{12}Când s-a apropiat de poarta cetății, iată că a fost dus afară un mort, singurul fiu născut al mamei sale, care era văduvă. Mulți oameni din cetate erau cu ea.

\VS{13}Domnul, când a văzut-o, i s-a făcut milă de ea și i-a zis:
{\WJ{“Nu plânge!”
}}

\VS{14}S-a apropiat și s-a atins de sicriu, iar purtătorii s-au oprit. El a spus:
{\WJ{“Tinere, îți spun: ridică-te!”.
}}

\VS{15}Cel care era mort s-a așezat și a început să vorbească. Apoi l-a dat mamei sale.


\par }{\PP \VS{16}Toți s-au temut și au slăvit pe Dumnezeu, zicând: “Un mare prooroc s-a ridicat printre noi!” și: “Dumnezeu a vizitat pe poporul Său!”.

\VS{17}Această veste s-a răspândit despre el în toată Iudeea și în toată regiunea înconjurătoare.


\par }{\PP \VS{18}Ucenicii lui Ioan i-au povestit toate acestea.

\VS{19}Ioan, chemând la el doi dintre ucenicii săi, i-a trimis la Isus, zicând: “Tu ești Cel ce vine, sau trebuie să căutăm pe altul?”

\VS{20}După ce au venit la el, au spus: “Ioan Botezătorul ne-a trimis la tine, spunând: “Tu ești cel care vine, sau trebuie să căutăm pe altul?””


\par }{\PP \VS{21}În ceasul acela a vindecat pe mulți de boli, de molime și de duhuri rele, și multora dintre cei orbi le-a dat vederea.

\VS{22}Isus le-a răspuns:
{\WJ{“Duceți-vă și spuneți lui Ioan lucrurile pe care le-ați văzut și le-ați auzit: că orbii își primesc vederea, șchiopii umblă, leproșii sunt curățiți, surzii aud, morții înviază și săracilor li se vestește o veste bună.
}}

\VS{23}{\WJ{Ferice de cel care nu găsește în mine nici un prilej de poticnire.”
}}


\par }{\PP \VS{24}După ce au plecat trimișii lui Ioan, a început să spună mulțimii despre Ioan: “{\WJ{Ce ați ieșit în pustie ca să vedeți? O trestie scuturată de vânt?
}}

\VS{25}{\WJ{Dar voi ce ați ieșit să vedeți? Un om îmbrăcat în haine moi? Iată, cei care sunt îmbrăcați frumos și trăiesc delicat sunt la curțile regilor.
}}

\VS{26}{\WJ{Dar tu ce ai ieșit să vezi? Un profet? Da, vă spun, și chiar mai mult decât un profet.
}}

\VS{27}{\WJ{Acesta este cel despre care este scris,
}}

\par }{\Q {\WJ{“Iată, trimit pe mesagerul Meu înaintea feței tale,
}}

\par }{\QB {\WJ{care îți va pregăti calea înaintea ta.
}}


\par }{\PP \VS{28}“{\WJ{Căci vă spun că, dintre cei născuți din femei, nu este un prooroc mai mare decât Ioan Botezătorul; dar cel mai mic în Împărăția lui Dumnezeu este mai mare decât el.”
}}


\par }{\PP \VS{29}Când a auzit aceasta, tot poporul și vameșii au declarat că Dumnezeu este drept, după ce a fost botezat cu botezul lui Ioan.

\VS{30}Dar fariseii și avocații au respins sfatul lui Dumnezeu, nefiind ei înșiși botezați de el.


\par }{\PP \VS{31}{\WJ{Cu ce aș putea să compar pe oamenii acestui neam? Cu ce se aseamănă ei?
}}

\VS{32}{\WJ{Sunt ca niște copii care stau în piață și se strigă unii pe alții, zicând: “Noi v-am cântat la cimpoi, și voi n-ați dansat”. Noi am jelit, iar voi nu ați plâns'.
}}

\VS{33}{\WJ{Căci Ioan Botezătorul n-a venit nici mâncând pâine, nici bând vin, iar voi spuneți: 'Are un demon'.
}}

\VS{34}{\WJ{Fiul Omului a venit mâncând și bând și voi ziceți: 'Iată un mâncăcios și un bețiv, un prieten al vameșilor și al păcătoșilor!'.
}}

\VS{35}{\WJ{Înțelepciunea este îndreptățită de toți copiii ei.”
}}


\par }{\PP \VS{36}Unul dintre farisei l-a invitat să mănânce cu el. El a intrat în casa fariseului și s-a așezat la masă.

\VS{37}Iată că o femeie din cetate, care era păcătoasă, când a aflat că el se odihnește în casa fariseului, a adus un vas de alabastru cu unguent.

\VS{38}Stând în spate la picioarele lui, plângând, a început să îi ude picioarele cu lacrimile ei, le-a șters cu părul capului, i-a sărutat picioarele și le-a uns cu mirul.

\VS{39}Fariseul care îl invitase, văzând aceasta, și-a zis în sinea lui: “Omul acesta, dacă ar fi fost profet, și-ar fi dat seama cine și ce fel de femeie este cea care îl atinge, că este o păcătoasă.”


\par }{\PP \VS{40}Isus i-a răspuns:
{\WJ{“Simon, am să-ți spun ceva.”
}}

\par }{\PP El a spus: “Învățătorule, spune mai departe”.


\par }{\PP \VS{41}“Un
{\WJ{creditor avea doi datornici. Unul datora cinci sute de denari, iar celălalt cincizeci.
}}

\VS{42}{\WJ{Când n-au putut plăti, i-a iertat pe amândoi. Așadar, care dintre ei îl va iubi mai mult?”
}}


\par }{\PP \VS{43}Simon a răspuns: “Cred că pe acela pe care l-a iertat cel mai mult.”

\par }{\PP El i-a spus:
{\WJ{“Ai judecat corect”.
}}

\VS{44}Și, întorcându-se spre femeie, i-a zis lui Simon: “{\WJ{Vezi tu pe femeia aceasta? Am intrat în casa ta și nu mi-ai dat apă pentru picioarele mele, dar ea mi-a udat picioarele cu lacrimile ei și le-a șters cu părul capului ei.
}}

\VS{45}{\WJ{Tu nu mi-ai dat nici o sărutare, dar ea, de când am intrat, nu încetează să-mi sărute picioarele.
}}

\VS{46}{\WJ{Tu nu mi-ai uns capul cu untdelemn, dar ea mi-a uns picioarele cu mir.
}}

\VS{47}{\WJ{De aceea vă spun că păcatele ei, care sunt multe, sunt iertate, pentru că a iubit mult. Dar cel căruia i se iartă puțin, iubește puțin.”
}}

\VS{48}El i-a zis:
{\WJ{“Păcatele tale sunt iertate”.
}}


\par }{\PP \VS{49}Și cei ce ședeau la masă cu El au început să se întrebe: “Cine este Acesta, care chiar iartă păcatele?”


\par }{\PP \VS{50}El a zis femeii: “{\WJ{Credința ta te-a mântuit. Du-te în pace.”
}}



\par }\Chap{8}{\PP \VerseOne{1}Curând după aceea, a străbătut cetățile și satele, propovăduind și aducând vestea cea bună a Împărăției lui Dumnezeu. Împreună cu el erau cei doisprezece

\VS{2}și câteva femei care fuseseră vindecate de duhuri rele și de boli: Maria, care se numea Magdalena, din care ieșiseră șapte demoni,

\VS{3}Ioana, soția lui Chuzas, administratorul lui Irod, Susana și multe altele care îi slujeau din averile lor.

\VS{4}Când s-a adunat o mare mulțime și veneau la el oameni din toate cetățile, el a vorbit printr-o pildă:

\VS{5}{\WJ{“Fermierul a ieșit să-și semene sămânța. Pe când semăna, o parte a căzut pe drum, a fost călcată în picioare și păsările cerului au mâncat-o.
}}

\VS{6}{\WJ{Altă sămânță a căzut pe stâncă și, de îndată ce a crescut, s-a uscat, pentru că nu avea umiditate.
}}

\VS{7}{\WJ{Altele au căzut în mijlocul spinilor, iar spinii au crescut cu ea și au înăbușit-o.
}}

\VS{8}{\WJ{Altă a căzut în pământul bun, a crescut și a produs de o sută de ori mai mult rod.” În
}}timp ce spunea aceste lucruri, a strigat:
{\WJ{“Cine are urechi de auzit, să audă!”
}}


\par }{\PP \VS{9}Atunci ucenicii Lui L-au întrebat: “Ce înseamnă pilda aceasta?”


\par }{\PP \VS{10}El a zis: “{\WJ{Vouă v-a fost dat să cunoașteți tainele Împărăției lui Dumnezeu; dar celorlalți le-a fost dat în pilde, ca să nu vadă, văzând, și să nu înțeleagă, auzind.
}}


\par }{\PP \VS{11}{\WJ{“Și pilda este următoarea: Sămânța este cuvântul lui Dumnezeu.
}}

\VS{12}Cei de
{\WJ{pe drum sunt cei care ascultă; apoi vine diavolul și le ia cuvântul din inimă, ca să nu creadă și să nu fie mântuiți.
}}

\VS{13}{\WJ{Cei de pe stâncă sunt cei care, când aud, primesc cuvântul cu bucurie; dar aceștia nu au rădăcină. Ei cred pentru o vreme, apoi cad în timp de ispită.
}}

\VS{14}{\WJ{Ceea ce a căzut între spini, aceștia sunt cei care au auzit și, mergând pe drumul lor, sunt sufocați de griji, de bogății și de plăcerile vieții; și nu aduc niciun rod la maturitate.
}}

\VS{15}{\WJ{Cele din pământul bun, aceștia sunt cei care, cu inimă cinstită și bună, după ce au auzit cuvântul, îl țin strâns și produc roade cu perseverență.
}}


\par }{\PP \VS{16}Nimeni, când
{\WJ{aprinde o lampă, nu o acoperă cu un vas și nu o pune sub pat, ci o pune pe un suport, ca să vadă lumina cei ce intră înăuntru.
}}

\VS{17}{\WJ{Căci nu este nimic ascuns care să nu fie descoperit, nici nimic secret care să nu fie cunoscut și să nu iasă la lumină.
}}

\VS{18}{\WJ{Fiți deci atenți cum auziți. Căci oricui are, i se va da; și oricui nu are, i se va lua chiar și ceea ce crede că are.”
}}


\par }{\PP \VS{19}Mama și frații lui au venit la el, dar nu puteau să se apropie de el din cauza mulțimii.

\VS{20}Unii oameni i-au spus: “Mama ta și frații tăi stau afară și doresc să te vadă.”


\par }{\PP \VS{21}Dar El le-a răspuns: “{\WJ{Mama mea și frații mei sunt aceia care au auzit cuvântul lui Dumnezeu și-l împlinesc.”
}}


\par }{\PP \VS{22}Într-una din zilele acelea, S-a suit într-o barcă, El însuși și ucenicii Săi, și le-a zis:
{\WJ{“Să trecem de cealaltă parte a lacului.” Așa că au pornit la apă.
}}

\VS{23}Dar, în timp ce navigau, el a adormit. O furtună de vânt s-a abătut asupra lacului, iar ei luau cantități periculoase de apă.

\VS{24}Au venit la el și l-au trezit, spunând: “Stăpâne, stăpâne, murim!” El s-a trezit și a mustrat vântul și furia apei; apoi au încetat și s-a făcut liniște.

\VS{25}El i-a întrebat:
{\WJ{“Unde este credința voastră?” Înfricoșați, ei s-au mirat, spunându-și unul altuia: “Cine este acesta, care poruncește chiar și vânturilor și apelor, iar acestea îl ascultă?”
}}


\par }{\PP \VS{26}Apoi au ajuns în ținutul Gadarenilor, care este în fața Galileii.

\VS{27}Când Isus a coborât la țărm, L-a întâmpinat un om din cetate, care avea demoni de multă vreme. Nu purta haine și nu locuia într-o casă, ci în morminte.

\VS{28}Când L-a văzut pe Isus, a strigat, a căzut înaintea Lui și a zis cu glas tare: “Ce am eu de-a face cu Tine, Isus, Fiul Dumnezeului Celui Preaînalt? Te implor, nu mă chinui!”.

\VS{29}Căci Isus poruncea duhului necurat să iasă din omul acela. Căci duhul necurat îl apucase de multe ori pe om. El era ținut sub pază și legat cu lanțuri și lanțuri. Rupând legăturile, demonul l-a împins în deșert.


\par }{\PP \VS{30}Isus l-a întrebat:
{\WJ{“Care este numele tău?”
}}

\par }{\PP El a spus: “Legiune”, pentru că mulți demoni intraseră în el.

\VS{31}Ei l-au implorat să nu le poruncească să se ducă în abis.


\par }{\PP \VS{32}Pe munte era o turmă de mulți porci care pășteau pe munte, și l-au rugat să le dea voie să intre în ei. Atunci el le-a îngăduit.

\VS{33}Demonii au ieșit din omul acela și au intrat în porci, iar turma s-a repezit pe malul abrupt în lac și s-a înecat.


\par }{\PP \VS{34}Când au văzut cei ce-i hrăneau ce se întâmplase, au fugit și au povestit în cetate și în țară.


\par }{\PP \VS{35}Oamenii au ieșit să vadă ce s-a întâmplat. Au venit la Isus și l-au găsit pe omul din care ieșiseră demonii, șezând la picioarele lui Isus, îmbrăcat și în deplinătatea facultăților mintale; și s-au temut.

\VS{36}Cei care l-au văzut le-au povestit cum s-a vindecat cel care fusese posedat de demoni.

\VS{37}Toți locuitorii din ținutul Gadarenilor din împrejurimi l-au rugat să se îndepărteze de ei, pentru că se temeau foarte mult. Atunci a intrat în barcă și s-a întors.

\VS{38}Omul din care ieșiseră demonii l-a implorat să meargă cu el, dar Isus l-a alungat, spunându-i:

\VS{39}{\WJ{“Întoarce-te la casa ta și povestește ce lucruri mari a făcut Dumnezeu pentru tine.”}} El a plecat, proclamând în toată cetatea ce lucruri mari a făcut Isus pentru el.


\par }{\PP \VS{40}Când s-a întors Isus, mulțimea L-a întâmpinat, căci toți Îl așteptau.

\VS{41}Și iată că a venit un om cu numele Iair. Era un conducător al sinagogii. A căzut la picioarele lui Isus și l-a rugat să intre în casa lui,

\VS{42}pentru că avea o fiică unică, în vârstă de vreo doisprezece ani, și era pe moarte. Dar, pe când se ducea, mulțimea îl apăsa.

\VS{43}O femeie care avea o scurgere de sânge de doisprezece ani, care cheltuise tot ce trăia pe medici și nu putea fi vindecată de niciunul,

\VS{44}a venit în spatele lui și s-a atins de franjurii hainei lui. Imediat curgerea sângelui ei s-a oprit.


\par }{\PP \VS{45}Isus a zis:
{\WJ{“Cine M-a atins?”
}}

\par }{\PP Când toți au negat, Petru și cei care erau cu el au zis: “Învățătorule, mulțimile te presează și te îmbrâncesc, iar Tu spui:
{\WJ{“Cine s-a atins de Mine?””.
}}


\par }{\PP \VS{46}Dar Isus a zis: “{\WJ{Cineva m-a atins, căci am văzut că a ieșit din mine o putere.”
}}

\VS{47}Când femeia a văzut că nu era ascunsă, a venit tremurând; și, căzând înaintea lui, i-a declarat în fața întregului popor motivul pentru care se atinsese de el și cum s-a vindecat imediat.

\VS{48}El i-a zis:
{\WJ{“Fiică, înveselește-te! Credința ta te-a făcut sănătoasă. Du-te în pace”.
}}


\par }{\PP \VS{49}Pe când vorbea el încă, a venit unul din casa conducătorului sinagogii și i-a zis: “Fiica ta a murit. Nu-l deranja pe Învățător”.


\par }{\PP \VS{50}Dar Isus, auzind, i-a răspuns: “{\WJ{Nu te teme. Crede numai și ea va fi vindecată.”
}}


\par }{\PP \VS{51}Când a ajuns la casă, n-a lăsat pe nimeni să intre înăuntru, decât pe Petru, pe Ioan, pe Iacov, pe tatăl fetiței și pe mama ei.

\VS{52}Toți plângeau și o jeleau, dar el a zis:
{\WJ{“Nu plângeți. Ea nu este moartă, ci doarme”.
}}


\par }{\PP \VS{53}Îl batjocoreau, știind că ea era moartă.

\VS{54}Dar el i-a scos pe toți afară și, luând-o de mână, a strigat-o și a zis:
{\WJ{“Copila, scoală-te!”
}}

\VS{55}Spiritul ei s-a întors și s-a sculat îndată. El a poruncit să i se dea ceva de mâncare.

\VS{56}Părinții ei au fost uimiți, dar el le-a poruncit să nu spună nimănui ce s-a întâmplat.



\par }\Chap{9}{\PP \VerseOne{1}A chemat laolaltă pe cei doisprezece și le-a dat putere și stăpânire asupra tuturor demonilor și să vindece bolile.

\VS{2}I-a trimis să propovăduiască Împărăția lui Dumnezeu și să vindece bolnavii.

\VS{3}El le-a spus:
{\WJ{“Nu luați nimic pentru călătoria voastră — nici toiag, nici portofel, nici pâine, nici bani. Nu aveți câte două tunici fiecare.
}}

\VS{4}{\WJ{În orice casă veți intra, rămâneți acolo și plecați de acolo.
}}

\VS{5}{\WJ{Toți cei care nu vă primesc, când veți pleca din cetatea aceea, scuturați-vă până și praful de pe picioare, ca mărturie împotriva lor.”
}}


\par }{\PP \VS{6}Și au plecat și au străbătut satele, propovăduind vestea cea bună și vindecând pretutindeni.


\par }{\PP \VS{7}Iar Irod tetrarhul a auzit tot ce făcuse el și a fost foarte tulburat, pentru că unii spuneau că Ioan a înviat din morți,

\VS{8}alții că Ilie s-a arătat, iar alții că a înviat unul din vechii prooroci.

\VS{9}Irod a spus: “Eu l-am decapitat pe Ioan, dar cine este acesta despre care aud astfel de lucruri?”. A căutat să îl vadă.


\par }{\PP \VS{10}Apostolii, după ce s-au întors, I-au povestit ce făcuseră.

\par }{\PP I-a luat și s-a retras deoparte, într-o regiune pustie, într-o cetate numită Betsaida.

\VS{11}Dar mulțimile, dându-și seama, l-au urmat. El i-a întâmpinat, le-a vorbit despre Împărăția lui Dumnezeu și i-a vindecat pe cei care aveau nevoie de vindecare.

\VS{12}Ziua a început să se scurgă, iar cei doisprezece au venit și i-au spus: “Trimiteți mulțimea să se ducă în satele și fermele din jur, să se cazeze și să facă rost de hrană, pentru că noi suntem aici într-un loc pustiu.”


\par }{\PP \VS{13}Dar El le-a zis: “{\WJ{Dați-le voi să mănânce ceva.”
}}

\par }{\PP Ei au spus: “Nu avem mai mult de cinci pâini și doi pești, dacă nu cumpărăm mâncare pentru toți acești oameni.”

\VS{14}Căci erau cam cinci mii de oameni.

\par }{\PP El le-a spus ucenicilor săi:
{\WJ{“Fă-i să se așeze în grupuri de câte cincizeci de persoane.”
}}

\VS{15}Ei au făcut așa și i-au pus pe toți să se așeze.

\VS{16}El a luat cele cinci pâini și cei doi pești și, privind spre cer, le-a binecuvântat, le-a frânt și le-a dat discipolilor ca să le pună înaintea mulțimii.

\VS{17}Ei au mâncat și s-au săturat cu toții. Au adunat douăsprezece coșuri cu bucățile sparte care rămăseseră.


\par }{\PP \VS{18}Pe când se ruga El singur, ucenicii erau lângă El și El i-a întrebat: “{\WJ{Cine zice mulțimea că sunt Eu?”
}}


\par }{\PP \VS{19}Ei au răspuns: “Ioan Botezătorul”; dar alții zic: “Ilie”, iar alții zic că a înviat unul din vechii prooroci”.


\par }{\PP \VS{20}El le-a zis: “{\WJ{Dar voi cine ziceți că sunt Eu?”
}}

\par }{\PP Petru a răspuns: “Hristosul lui Dumnezeu.”


\par }{\PP \VS{21}Dar El i-a avertizat și le-a poruncit să nu spună nimănui acest lucru,

\VS{22}zicând: “{\WJ{Fiul Omului trebuie să sufere multe, să fie lepădat de bătrâni, de preoții cei mai de seamă și de cărturari, să fie omorât și a treia zi să învieze.”
}}


\par }{\PP \VS{23}Și a zis tuturor: “{\WJ{Dacă voiește cineva să vină după Mine, să se lepede de sine, să-și ia crucea și să-Mi urmeze Mie.
}}

\VS{24}{\WJ{Căci oricine vrea să-și salveze viața o va pierde, dar oricine își va pierde viața pentru Mine o va salva.
}}

\VS{25}{\WJ{Căci ce-i folosește cuiva să câștige toată lumea și să se piardă sau să se piardă pe sine însuși?
}}

\VS{26}{\WJ{Căci oricine se va rușina de Mine și de cuvintele Mele, de acela se va rușina Fiul Omului când va veni în slava Sa și în slava Tatălui și a sfinților îngeri.
}}

\VS{27}{\WJ{Dar vă spun adevărul: sunt unii dintre cei care stau aici care nu vor gusta nicidecum moartea până când nu vor vedea Împărăția lui Dumnezeu.”
}}


\par }{\PP \VS{28}La vreo opt zile după aceste cuvinte, a luat cu Sine pe Petru, pe Ioan și pe Iacov și s-a suit pe munte să se roage.

\VS{29}În timp ce se ruga, înfățișarea feței Lui s-a schimbat, iar hainele Lui au devenit albe și orbitoare.

\VS{30}Iată că vorbeau cu el doi bărbați, care erau Moise și Ilie,

\VS{31}care s-au arătat în slavă și i-au vorbit despre plecarea pe care urma să o îndeplinească la Ierusalim.


\par }{\PP \VS{32}Petru și cei ce erau cu el erau îngreuiați de somn; dar, când s-au trezit, au văzut slava Lui și pe cei doi bărbați care stăteau cu El.

\VS{33}Pe când se despărțeau de El, Petru i-a zis lui Isus: “Învățătorule, este bine să fim aici. Să facem trei corturi: unul pentru tine, unul pentru Moise și unul pentru Ilie”, fără să știe ce spune.


\par }{\PP \VS{34}Pe când zicea El aceste lucruri, un nor a venit și i-a acoperit și ei s-au temut, când au intrat în nor.

\VS{35}Un glas a ieșit din nor și a zis: “Acesta este Fiul Meu preaiubit. Ascultați-l!”

\VS{36}Când a venit vocea, Isus a fost găsit singur. Ei au tăcut și nu au povestit nimănui, în acele zile, nimic din ceea ce văzuseră.


\par }{\PP \VS{37}A doua zi, când s-au coborât de pe munte, o mare mulțime de oameni I-a ieșit în întâmpinare.

\VS{38}Și iată că un om din mulțime a strigat: “Învățătorule, te rog să te uiți la fiul meu, căci este singurul meu copil născut.

\VS{39}Iată că un duh îl ia, strigă deodată și îl convulsionează până la spumegarea lui; și abia se depărtează de el, lovindu-l grav.

\VS{40}I-am rugat pe discipolii tăi să-l alunge, dar nu au putut.”


\par }{\PP \VS{41}Isus a răspuns: “{\WJ{Neam necredincios și pervers, până când voi fi cu voi și voi răbda cu voi? Aduceți aici pe fiul vostru”.
}}


\par }{\PP \VS{42}Pe când venea el încă, demonul l-a aruncat la pământ și l-a zdrobit cu putere. Dar Isus a mustrat duhul necurat, a vindecat băiatul și l-a dat înapoi tatălui său.

\VS{43}Toți erau uimiți de măreția lui Dumnezeu.

\par }{\PP Și, pe când toți se mirau de toate lucrurile pe care le făcea Isus, a zis ucenicilor Săi:

\VS{44}{\WJ{“Lăsați să vă intre în urechi cuvintele acestea, căci Fiul Omului va fi dat în mâinile oamenilor”.
}}

\VS{45}Dar ei nu au înțeles acest cuvânt. Le era ascunsă, ca să nu o înțeleagă, și le era teamă să îl întrebe despre acest cuvânt.


\par }{\PP \VS{46}Și s-a iscat între ei o ceartă despre care dintre ei era cel mai mare.

\VS{47}Isus, pricepând rațiunile inimilor lor, a luat un copilaș, l-a pus lângă El și le-a spus:

\VS{48}{\WJ{“Oricine primește pe acest copilaș în numele Meu, pe Mine Mă primește. Oricine mă primește pe mine îl primește pe cel care m-a trimis. Căci oricine este cel mai mic dintre voi toți, acesta va fi mare”.
}}


\par }{\PP \VS{49}Ioan a răspuns: “Învățătorule, am văzut pe cineva care scotea demoni în Numele Tău și i-am interzis, pentru că nu este de partea noastră.”


\par }{\PP \VS{50}Isus i-a zis:
{\WJ{“Nu-l opri, căci cine nu este împotriva noastră, este pentru noi.”
}}


\par }{\PP \VS{51}Când se apropiau zilele în care trebuia să fie luat, el și-a pus în gând să meargă la Ierusalim

\VS{52}și a trimis mesageri înaintea lui. Aceștia s-au dus și au intrat într-un sat al samaritenilor, ca să se pregătească pentru el.

\VS{53}Ei nu l-au primit, pentru că el călătorea cu fața îndreptată spre Ierusalim.

\VS{54}Când discipolii săi, Iacov și Ioan, au văzut acest lucru, au zis: “Doamne, vrei să poruncim noi să coboare foc din cer și să-i nimicească, așa cum a făcut Ilie?”


\par }{\PP \VS{55}Dar El s-a întors și i-a mustrat: “{\WJ{Nu știți ce fel de duh sunteți.
}}

\VS{56}{\WJ{Căci Fiul Omului nu a venit să distrugă viețile oamenilor, ci să le salveze.”
}}

\par }{\PP S-au dus în alt sat.

\VS{57}Pe când mergeau pe drum, un om I-a zis: “Vreau să Te urmez oriunde vei merge, Doamne.”


\par }{\PP \VS{58}Isus i-a zis: “{\WJ{Vulpile au vizuini și păsările cerului au cuiburi, dar Fiul Omului nu are unde să-Și pună capul.”
}}


\par }{\PP \VS{59}Și a zis altuia:
{\WJ{“Urmează-mă!”
}}

\par }{\PP Dar el a spus: “Doamne, dă-mi voie să mă duc mai întâi să-l îngrop pe tatăl meu”.


\par }{\PP \VS{60}Dar Isus i-a zis: “{\WJ{Lasă morții să-și îngroape morții lor, iar tu du-te și vestește Împărăția lui Dumnezeu.”
}}


\par }{\PP \VS{61}Un altul a mai zis: “Vreau să te urmez, Doamne, dar mai întâi îngăduie-mi să-mi iau rămas bun de la cei care sunt la mine acasă.”


\par }{\PP \VS{62}Dar Isus i-a zis: “{\WJ{Nimeni care a pus mâna la plug și se uită înapoi nu este potrivit pentru Împărăția lui Dumnezeu.”
}}



\par }\Chap{10}{\PP \VerseOne{1}După aceste lucruri, Domnul a rânduit și pe alți șaptezeci și i-a trimis înaintea Lui, doi câte doi, în fiecare cetate și în fiecare loc unde avea să vină.

\VS{2}Și le-a spus: “{\WJ{Secerișul este într-adevăr mare, dar lucrătorii sunt puțini. Rugați-vă, așadar, Domnului secerișului, ca să trimită lucrători la secerișul Său.
}}

\VS{3}{\WJ{Mergeți și voi pe căile voastre. Iată, vă trimit ca pe niște miei în mijlocul lupilor.
}}

\VS{4}{\WJ{Nu purtați nici pungă, nici portofel, nici sandale. Nu salutați pe nimeni pe drum.
}}

\VS{5}{\WJ{În orice casă veți intra, spuneți mai întâi: “Pacea fie cu casa aceasta!”.
}}

\VS{6}{\WJ{Dacă este acolo un fiu al păcii, pacea voastră se va odihni asupra lui; dacă nu, se va întoarce la voi.
}}

\VS{7}{\WJ{Rămâi în aceeași casă, mâncând și bând ce ți se dă, căci lucrătorul este vrednic de plata sa. Nu mergeți din casă în casă.
}}

\VS{8}{\WJ{În orice cetate veți intra și vă vor primi, mâncați ceea ce vi se pune înainte.
}}

\VS{9}{\WJ{Vindecați pe bolnavii care se află acolo și spuneți-le: “Împărăția lui Dumnezeu s-a apropiat de voi”.
}}

\VS{10}{\WJ{Dar în orice cetate în care intrați și nu vă primesc, ieșiți pe străzile ei și spuneți:
}}

\VS{11}“{\WJ{Chiar și praful din cetatea voastră care se lipește de noi, îl ștergem împotriva voastră. Cu toate acestea, să știți că Împărăția lui Dumnezeu s-a apropiat de voi'.
}}

\VS{12}{\WJ{Vă spun că, în ziua aceea, va fi mai ușor de suportat pentru Sodoma decât pentru orașul acela.
}}


\par }{\PP \VS{13}{\WJ{“Vai de tine, Corazin! Vai de tine, Betsaida! Căci dacă s-ar fi făcut în Tir și în Sidon lucrările mărețe care s-au făcut în voi, s-ar fi pocăit de mult, stând în sac și cenușă.
}}

\VS{14}{\WJ{Dar, la judecată, Tirului și Sidonului le va fi mai ușor de suportat decât vouă.
}}

\VS{15}{\WJ{Tu, Capernaum, care te-ai înălțat la cer, vei fi coborât în Hades.
}}

\VS{16}{\WJ{Cine vă ascultă pe voi mă ascultă pe mine, iar cine vă respinge pe voi mă respinge pe mine. Cine mă respinge pe mine, îl respinge pe cel care m-a trimis pe mine.”
}}


\par }{\PP \VS{17}Cei șaptezeci s-au întors cu bucurie, zicând: “Doamne, chiar și demonii ne sunt supuși în Numele Tău!”


\par }{\PP \VS{18}Și le-a zis: “{\WJ{Am văzut pe Satana căzând din cer ca un fulger.
}}

\VS{19}{\WJ{Iată, vă dau putere să călcați peste șerpi și scorpioni și peste toată puterea vrăjmașului. Nimic nu vă va răni în vreun fel.
}}

\VS{20}Cu toate acestea,
{\WJ{nu vă bucurați de faptul că duhurile vi se supun, ci bucurați-vă că numele voastre sunt scrise în ceruri.”
}}


\par }{\PP \VS{21}În același ceas, Isus s-a bucurat în Duhul Sfânt și a zis: “{\WJ{Îți mulțumesc, Tată, Stăpânul cerului și al pământului, că ai ascuns aceste lucruri de cei înțelepți și pricepuți și le-ai descoperit copiilor. Da, Tată, pentru că așa a fost bineplăcut în ochii tăi”.
}}


\par }{\PP \VS{22}Întorcându-Se către ucenici, a zis: “{\WJ{Toate Mi-au fost date de Tatăl Meu. Nimeni nu știe cine este Fiul, în afară de Tatăl, și cine este Tatăl, în afară de Fiul și de cel căruia Fiul vrea să i-l dezvăluie.”
}}


\par }{\PP \VS{23}Și, întorcându-se către ucenici, a zis în particular: “{\WJ{Fericiți ochii care văd ceea ce vedeți voi,
}}

\VS{24}{\WJ{căci vă spun că mulți prooroci și împărați au dorit să vadă ceea ce vedeți voi și n-au văzut, și să audă ceea ce auziți voi și n-au auzit.”
}}


\par }{\PP \VS{25}Și iată că un învățător s-a ridicat în picioare și L-a pus la încercare, zicând: “Învățătorule, ce să fac ca să moștenesc viața veșnică?”


\par }{\PP \VS{26}El i-a zis: “{\WJ{Ce este scris în Lege? Cum o citești?”
}}


\par }{\PP \VS{27}El a răspuns: “Să iubești pe Domnul Dumnezeul tău din toată inima ta, din tot sufletul tău, din toată puterea ta și din tot cugetul tău, și pe aproapele tău ca pe tine însuți.”


\par }{\PP \VS{28}El i-a zis: “Ai
{\WJ{răspuns corect. Fă așa și vei trăi.”
}}


\par }{\PP \VS{29}Dar el, vrând să se justifice, a întrebat pe Isus: “Cine este aproapele meu?”


\par }{\PP \VS{30}Isus a răspuns: “{\WJ{Un om care se cobora de la Ierusalim la Ierihon, a căzut în mijlocul unor tâlhari care, după ce l-au dezbrăcat și l-au bătut, au plecat și l-au lăsat pe jumătate mort.
}}

\VS{31}Din
{\WJ{întâmplare, un preot cobora pe drumul acela. Când l-a văzut, a trecut pe partea cealaltă.
}}

\VS{32}La
{\WJ{fel și un levit, când a ajuns la locul acela și l-a văzut, a trecut pe partea cealaltă.
}}

\VS{33}{\WJ{Dar un oarecare samaritean, pe când călătorea, a ajuns unde era el. Când l-a văzut, i s-a făcut milă,
}}

\VS{34}s-a
{\WJ{apropiat de el și i-a legat rănile, turnând untdelemn și vin. L-a așezat pe animalul său, l-a dus la un han și a avut grijă de el.
}}

\VS{35}{\WJ{A doua zi, când a plecat, a scos doi denari, i-a dat gazdei și i-a spus: “Ai grijă de el. Tot ce vei cheltui în plus, îți voi da înapoi când mă voi întoarce'.
}}

\VS{36}{\WJ{Acum, care dintre acești trei credeți că vi s-a părut că era vecin cu cel care a căzut între tâlhari?”
}}


\par }{\PP \VS{37}El a zis: “Cel care a avut milă de el.”

\par }{\PP Atunci Isus i-a zis:
{\WJ{“Du-te și fă și tu la fel”.
}}


\par }{\PP \VS{38}Pe când mergeau ei pe drum, a intrat într-un sat și o femeie, numită Marta, L-a primit în casa ei.

\VS{39}Ea avea o soră, numită Maria, care ședea și ea la picioarele lui Isus și asculta cuvântul Lui.

\VS{40}Dar Marta, care era ocupată cu multe slujbe, s-a apropiat de el și i-a zis: “Doamne, nu-ți pasă că sora mea m-a lăsat să slujesc singură? Roag-o, așadar, să mă ajute”.


\par }{\PP \VS{41}Isus i-a răspuns: “{\WJ{Marta, Marta, tu ești neliniștită și îngrijorată de multe lucruri;
}}

\VS{42}{\WJ{dar un singur lucru este necesar. Maria a ales partea cea bună, care nu-i va fi luată”.
}}



\par }\Chap{11}{\PP \VerseOne{1}După ce a terminat de rugat într-un loc, unul din ucenicii Săi I-a zis: “Doamne, învață-ne să ne rugăm, cum i-a învățat și Ioan pe ucenicii Săi.”


\par }{\PP \VS{2}El le-a zis: “{\WJ{Când vă rugați, spuneți,
}}

\par }{\Q {\WJ{“Tatăl nostru care ești în ceruri,
}}

\par }{\QB {\WJ{fie ca numele tău să fie păstrat sfânt.
}}

\par }{\Q {\WJ{Fie ca Împărăția Ta să vină.
}}

\par }{\QB {\WJ{Facă-se voia Ta, precum în cer așa și pe pământ.
}}


\par }{\Q \VS{3}{\WJ{Dă-ne în fiecare zi pâinea noastră cea de toate zilele.
}}


\par }{\Q \VS{4}{\WJ{Iartă-ne nouă păcatele noastre,
}}

\par }{\QB {\WJ{pentru că și noi înșine iertăm tuturor celor care ne sunt datori.
}}

\par }{\Q {\WJ{Nu ne duce pe noi în ispită,
}}

\par }{\QB {\WJ{ci izbăvește-ne pe noi de cel rău”.”
}}


\par }{\PP \VS{5}El le-a zis: “{\WJ{Care dintre voi, dacă vă veți duce la miezul nopții la un prieten și-i veți spune: “Prietene, împrumută-mi trei pâini,
}}

\VS{6}{\WJ{căci un prieten de-al meu a venit la mine din călătorie și nu am ce să-i pun înainte,
}}

\VS{7}{\WJ{iar el, dinăuntru, va răspunde și va zice: “Nu mă deranja. Ușa este acum închisă, iar copiii mei sunt cu mine în pat. Nu pot să mă ridic și să ți-i dau'?”?
}}

\VS{8}{\WJ{Vă spun că, deși nu se va ridica și nu i le va da, pentru că este prietenul lui, totuși, datorită insistenței lui, se va ridica și îi va da atâtea câte îi trebuie.
}}


\par }{\PP \VS{9}“{\WJ{Vă spun: cereți și vi se va da. Continuați să căutați și veți găsi. Continuați să bateți la ușă și vi se va deschide.
}}

\VS{10}{\WJ{Căci oricine cere primește. Cel care caută găsește. Celui care bate i se va deschide.
}}


\par }{\PP \VS{11}{\WJ{Care dintre voi, taților, dacă fiul vostru va cere o pâine, îi va da o piatră? Sau, dacă va cere un pește, nu-i va da un șarpe în loc de pește, nu-i așa?
}}

\VS{12}{\WJ{Sau, dacă va cere un ou, nu-i va da un scorpion, nu-i așa?
}}

\VS{13}{\WJ{Dacă voi, deci, fiind răi, știți să dați daruri bune copiilor voștri, cu cât mai mult Tatăl vostru ceresc va da Duhul Sfânt celor care i-l cer?”
}}


\par }{\PP \VS{14}El scotea un demon, și acesta era mut. După ce demonul a ieșit, muticul a vorbit; și mulțimile se mirau.

\VS{15}Dar unii dintre ei ziceau: “El scoate demonii prin Beelzebul, prințul demonilor.”

\VS{16}Alții, punându-l la încercare, căutau de la el un semn din cer.

\VS{17}Dar El, cunoscându-le gândurile, le-a zis: “{\WJ{Orice împărăție împărțită împotriva ei însăși este dărâmată. O casă împărțită împotriva ei însăși cade.
}}

\VS{18}{\WJ{Dacă și Satana este dezbinat împotriva lui însuși, cum va rezista împărăția lui? Căci voi spuneți că eu scot demoni prin Beelzebul.
}}

\VS{19}{\WJ{Dar dacă Eu scot demoni prin Beelzebul, prin cine îi scot copiii voștri? De aceea ei vor fi judecătorii voștri.
}}

\VS{20}{\WJ{Dar dacă eu, prin degetul lui Dumnezeu, scot afară demoni, atunci Împărăția lui Dumnezeu a venit la voi.
}}


\par }{\PP \VS{21}{\WJ{“Când un om puternic, înarmat până-n dinți, își păzește locuința, bunurile sale sunt în siguranță.
}}

\VS{22}{\WJ{Dar când cineva mai puternic îl atacă și îl învinge, îi ia toată armura în care se încredea și își împarte prada.
}}


\par }{\PP \VS{23}{\WJ{Cine nu este cu Mine, este împotriva Mea. Cine nu se adună cu mine se risipește.
}}


\par }{\PP \VS{24}După ce a ieșit
{\WJ{din om, duhul necurat trece prin locuri uscate, căutând odihnă; și, cum nu găsește nimic, zice: “Mă voi întoarce la casa mea de unde am ieșit.
}}

\VS{25}{\WJ{Când se întoarce, o găsește măturată și pusă în ordine.
}}

\VS{26}{\WJ{Atunci se duce și ia alte șapte duhuri mai rele decât el, intră și locuiesc acolo. Ultima stare a acelui om devine mai rea decât prima.”
}}


\par }{\PP \VS{27}Pe când zicea El aceste lucruri, o femeie din mulțime a ridicat glasul și I-a zis: “Binecuvântat este pântecele care Te-a născut și sânii care Te-au alăptat!”


\par }{\PP \VS{28}Iar el a zis: “Dimpotrivă,
{\WJ{fericiți cei ce ascultă cuvântul lui Dumnezeu și-l păzesc.”
}}


\par }{\PP \VS{29}Când s-a adunat mulțimea la El, a început să zică: “{\WJ{Acesta este un neam rău. Ea caută un semn. Nu i se va da alt semn decât semnul profetului Iona.
}}

\VS{30}{\WJ{Căci, după cum Iona a fost un semn pentru niniviteni, așa va fi și Fiul Omului pentru acest neam.
}}

\VS{31}{\WJ{Împărăteasa Sudului se va ridica la judecată împreună cu oamenii acestui neam și îi va condamna, pentru că a venit de la marginile pământului să asculte înțelepciunea lui Solomon; și iată că unul mai mare decât Solomon este aici.
}}

\VS{32}{\WJ{Bărbații din Ninive se vor ridica la judecată împreună cu acest neam și îl vor condamna, pentru că s-au pocăit la propovăduirea lui Iona; și iată, unul mai mare decât Iona este aici.
}}


\par }{\PP \VS{33}Nimeni, când
{\WJ{aprinde o lampă, nu o pune în pivniță sau sub un coș, ci pe un suport, ca să vadă lumina cei ce intră.
}}

\VS{34}{\WJ{Lampa trupului este ochiul. De aceea, când ochiul tău este bun, tot corpul tău este și el plin de lumină; dar când este rău, tot corpul tău este plin de întuneric.
}}

\VS{35}{\WJ{De aceea, vedeți dacă lumina care este în voi nu este întuneric.
}}

\VS{36}{\WJ{Așadar, dacă tot trupul vostru este plin de lumină, neavând nicio parte întunecată, va fi în întregime plin de lumină, ca atunci când lampa cu strălucirea ei strălucitoare vă luminează.”
}}


\par }{\PP \VS{37}Pe când vorbea el, un fariseu l-a invitat să ia masa cu el. El a intrat și s-a așezat la masă.

\VS{38}Fariseul, când l-a văzut, s-a mirat că nu se spălase mai întâi înainte de cină.

\VS{39}Domnul i-a zis: “{\WJ{Voi, fariseii, curățați partea exterioară a paharului și a farfuriei, dar partea voastră interioară este plină de șantaj și de răutate.
}}

\VS{40}{\WJ{Voi, nebunilor, nu cumva Cel care a făcut partea exterioară nu a făcut și partea interioară?
}}

\VS{41}{\WJ{Dar dați ca daruri celor nevoiași cele dinăuntru și iată că toate lucrurile vor fi curate pentru voi.
}}

\VS{42}{\WJ{Dar vai de voi, fariseilor! Pentru că voi dați zeciuială pentru mentă, rudenie și orice altă plantă, dar ocoliți dreptatea și dragostea lui Dumnezeu. Ar fi trebuit să le faceți pe acestea și să nu le lăsați pe celelalte nefăcute.
}}

\VS{43}{\WJ{Vai de voi, fariseilor! Pentru că vă plac locurile cele mai bune din sinagogi și salutul din piețe.
}}

\VS{44}{\WJ{Vai vouă, cărturari și farisei, ipocriților! Pentru că sunteți ca niște morminte ascunse, iar cei care umblă peste ele nu știu.”
}}


\par }{\PP \VS{45}Unul dintre avocați i-a răspuns: “Învățătorule, spunând aceasta, ne insulți și pe noi.”


\par }{\PP \VS{46}El a zis: “{\WJ{Vai de voi, avocaților! Pentru că îi încărcați pe oameni cu poveri greu de purtat, iar voi înșivă nu ridicați nici măcar un deget pentru a ajuta la purtarea acestor poveri.
}}

\VS{47}{\WJ{Vai de voi! Pentru că voi construiți mormintele profeților, iar părinții voștri i-au ucis.
}}

\VS{48}{\WJ{Așa că voi mărturisiți și consimțiți la faptele părinților voștri. Căci ei i-au ucis, iar voi le construiți mormintele.
}}

\VS{49}{\WJ{De aceea și înțelepciunea lui Dumnezeu a spus: “Voi trimite la ei profeți și apostoli; și pe unii dintre ei îi vor ucide și îi vor prigoni,
}}

\VS{50}pentru
{\WJ{ca sângele tuturor profeților, care a fost vărsat de la întemeierea lumii, să fie cerut de acest neam,
}}

\VS{51}de
{\WJ{la sângele lui Abel până la sângele lui Zaharia, care a pierit între altar și sanctuar. Da, vă spun, va fi cerut de această generație.
}}

\VS{52}{\WJ{Vai de voi, avocaților! Căci ați luat cheia cunoștinței. Voi înșivă nu ați intrat înăuntru, iar pe cei care intrau, voi i-ați împiedicat.”
}}


\par }{\PP \VS{53}Pe când le spunea El aceste lucruri, cărturarii și fariseii au început să se mânie cumplit și să scoată multe lucruri de la El,

\VS{54}pândindu-L și căutând să-L prindă în ceva ce ar putea să spună, ca să-L acuze.



\par }\Chap{12}{\PP \VerseOne{1}Între timp, după ce s-a adunat o mulțime de multe mii de oameni, încât se călcau în picioare unii pe alții, a început să spună mai întâi ucenicilor Săi: “Feriți-vă de
{\WJ{drojdia fariseilor, care este ipocrizia.
}}

\VS{2}{\WJ{Dar nu este nimic ascuns care să nu fie descoperit, nici ascuns care să nu fie cunoscut.
}}

\VS{3}{\WJ{De aceea, tot ce ați spus în întuneric va fi auzit la lumină. Ceea ce ați spus la ureche în încăperile interioare va fi proclamat pe acoperișurile caselor.
}}


\par }{\PP \VS{4}“{\WJ{Vă spun, prietenii mei, să nu vă temeți de cei ce ucid trupul, și după aceea nu mai au ce să facă.
}}

\VS{5}{\WJ{Ci vă voi avertiza de cine trebuie să vă temeți. Temeți-vă de cel care, după ce a ucis, are puterea de a arunca în Gheenă. Da, vă spun, temeți-vă de el.
}}


\par }{\PP \VS{6}{\WJ{“Nu se vând cinci vrăbii cu două monede de assaria? Nici una dintre ele nu este uitată de Dumnezeu.
}}

\VS{7}{\WJ{Dar până și firele de păr de pe capul vostru sunt toate numărate. De aceea nu vă temeți. Voi sunteți mai valoroși decât multe vrăbii.
}}


\par }{\PP \VS{8}“{\WJ{Vă spun că pe oricine Mă va mărturisi înaintea oamenilor, Fiul Omului îl va mărturisi și înaintea îngerilor lui Dumnezeu;
}}

\VS{9}{\WJ{iar cine Mă va tăgădui în fața oamenilor, va fi tăgăduit și în fața îngerilor lui Dumnezeu.
}}

\VS{10}{\WJ{Oricine va rosti un cuvânt împotriva Fiului Omului va fi iertat, dar cei care blasfemiază împotriva Duhului Sfânt nu vor fi iertați.
}}

\VS{11}{\WJ{Când vă vor aduce în fața sinagogilor, a conducătorilor și a autorităților, nu vă îngrijorați cum sau ce veți răspunde sau ce veți spune;
}}

\VS{12}{\WJ{pentru că Duhul Sfânt vă va învăța în același ceas ce trebuie să spuneți.”
}}


\par }{\PP \VS{13}Unul din mulțime I-a zis: “Învățătorule, spune fratelui meu să împartă cu mine moștenirea.”


\par }{\PP \VS{14}Iar el i-a zis: “{\WJ{Omule, cine m-a pus pe mine judecător sau arbitru peste tine?”
}}

\VS{15}El le-a zis:
{\WJ{“Luați seama! Păstrați-vă departe de lăcomie, căci viața unui om nu constă în abundența lucrurilor pe care le posedă”.
}}


\par }{\PP \VS{16}Și le-a spus o pildă, zicând: “{\WJ{Pământul unui om bogat producea din belșug.
}}

\VS{17}{\WJ{El se gândea în sinea lui și zicea: “Ce voi face, pentru că nu am loc să-mi depozitez recolta?”
}}

\VS{18}{\WJ{El a zis: “Iată ce voi face. Îmi voi dărâma hambarele, voi construi altele mai mari și acolo îmi voi depozita tot grâul și toate bunurile mele.
}}

\VS{19}{\WJ{Voi spune sufletului meu: “Suflete, ai multe bunuri depozitate de mulți ani. Ia-ți liniștea, mănâncă, bea și fii vesel!””'”.
}}


\par }{\PP \VS{20}{\WJ{Dar Dumnezeu i-a zis: “Nebunule, în seara aceasta ți se cere sufletul tău. Lucrurile pe care le-ai pregătit — ale cui vor fi?”
}}

\VS{21}{\WJ{Așa este cel care își strânge comori pentru sine și nu este bogat față de Dumnezeu.”
}}


\par }{\PP \VS{22}Și a zis ucenicilor Săi: “{\WJ{De aceea vă spun: Nu vă îngrijorați pentru viața voastră, ce veți mânca, nici pentru trupul vostru, ce veți purta.
}}

\VS{23}{\WJ{Viața este mai mult decât mâncarea, iar trupul este mai mult decât îmbrăcămintea.
}}

\VS{24}{\WJ{Gândiți-vă la corbi: ei nu seamănă, nu seceră, nu au depozit sau hambar, și Dumnezeu îi hrănește. Cât de mult mai valoroși sunteți voi decât păsările!
}}

\VS{25}{\WJ{Care dintre voi, fiind neliniștit, poate adăuga un cot la înălțimea sa?
}}

\VS{26}{\WJ{Dacă deci nu sunteți în stare să faceți nici măcar cele mai mici lucruri, de ce vă îngrijorați pentru restul?
}}

\VS{27}{\WJ{Gândiți-vă la crini, cum cresc ei. Ei nu se trudesc, nici nu se învârt; totuși vă spun că nici măcar Solomon, în toată slava lui, nu era îmbrăcat ca unul dintre aceștia.
}}

\VS{28}{\WJ{Dar dacă așa îmbracă Dumnezeu iarba de pe câmp, care astăzi există, iar mâine este aruncată în cuptor, cu cât mai mult vă va îmbrăca pe voi, cei cu puțină credință?
}}


\par }{\PP \VS{29}“{\WJ{Nu căutați ce veți mânca și ce veți bea, și nu vă îngrijorați.
}}

\VS{30}{\WJ{Căci neamurile lumii caută toate aceste lucruri, dar Tatăl vostru știe că voi aveți nevoie de ele.
}}

\VS{31}{\WJ{Ci căutați Împărăția lui Dumnezeu și toate aceste lucruri vi se vor adăuga.
}}


\par }{\PP \VS{32}“{\WJ{Nu vă temeți, turmă mică, căci Tatăl vostru a binevoit să vă dea Împărăția.
}}

\VS{33}{\WJ{Vindeți ce aveți și dați daruri celor nevoiași. Făceți-vă poșete care nu îmbătrânesc, o comoară în ceruri care nu se strică, de care nu se apropie niciun hoț și pe care nu o distruge molia.
}}

\VS{34}{\WJ{Căci unde este comoara voastră, acolo va fi și inima voastră.
}}


\par }{\PP \VS{35}“Să vă
{\WJ{îmbrăcați la brâu și să vă aprindeți lămpile.
}}

\VS{36}Să
{\WJ{fiți ca niște oameni care pândesc pe stăpânul lor când se întoarce de la nuntă, pentru ca, atunci când vine și bate la ușă, să-i deschidă îndată.
}}

\VS{37}{\WJ{Fericiți sunt acei slujitori pe care Domnul îi va găsi veghind când va veni. Mai mult ca sigur vă spun că se va îmbrăca, îi va face să se culce și va veni să le slujească.
}}

\VS{38}{\WJ{Ei vor fi binecuvântați dacă va veni în a doua sau a treia gardă și îi va găsi așa.
}}

\VS{39}{\WJ{Dar să știți că, dacă stăpânul casei ar fi știut la ce oră va veni hoțul, ar fi vegheat și nu ar fi permis să i se spargă casa.
}}

\VS{40}{\WJ{De aceea și voi fiți gata, căci Fiul Omului vine într-un ceas în care nu vă așteptați.”
}}


\par }{\PP \VS{41}Petru I-a zis: “Doamne, ne spui nouă pilda aceasta, sau tuturor?”


\par }{\PP \VS{42}Și Domnul a zis: “{\WJ{Cine este administratorul credincios și înțelept, pe care îl va pune stăpânul său peste casa lui, ca să le dea partea lor de mâncare la vremea potrivită?
}}

\VS{43}{\WJ{Ferice de acel slujitor pe care stăpânul său îl va găsi făcând așa când va veni.
}}

\VS{44}{\WJ{Adevărat vă spun că îl va pune peste tot ce are.
}}

\VS{45}{\WJ{Dar dacă robul acela zice în inima lui: “Domnul meu întârzie să vină” și începe să bată pe slujitorii și pe slujnice, să mănânce, să bea și să se îmbete,
}}

\VS{46}{\WJ{atunci stăpânul robului acela va veni într-o zi în care nu-l așteaptă și într-un ceas pe care nu-l știe, îl va tăia în două și-i va pune partea lui cu cei necredincioși.
}}

\VS{47}{\WJ{Robul acela care știa voia stăpânului său și nu s-a pregătit și nici nu a făcut ce voia el, va fi bătut cu multe lovituri,
}}

\VS{48}{\WJ{dar cel care nu știa și a făcut lucruri vrednice de lovituri va fi bătut cu puține lovituri. Cui i s-a dat mult, i se va cere mult; și cui i s-a încredințat mult, i se va cere mai mult.
}}


\par }{\PP \VS{49}{\WJ{“Am venit să arunc foc pe pământ. Aș vrea să fie deja aprins.
}}

\VS{50}{\WJ{Dar eu am un botez cu care trebuie să fiu botezat și cât de mult mă chinui până când se va împlini!
}}

\VS{51}{\WJ{Credeți că am venit să dau pace pe pământ? Vă spun că nu, ci mai degrabă dezbinare.
}}

\VS{52}{\WJ{Căci, de acum înainte, într-o casă vor fi cinci dezbinați, trei împotriva a doi și doi împotriva a trei.
}}

\VS{53}{\WJ{Vor fi dezbinați: tatăl împotriva fiului și fiul împotriva tatălui; mama împotriva fiicei și fiica împotriva mamei sale; soacra împotriva nurorii și nora împotriva soacrei.”
}}


\par }{\PP \VS{54}Și a zis mulțimilor: “{\WJ{Când vedeți un nor ridicându-se dinspre apus, îndată ziceți: “Vine o ploaie”, și așa se întâmplă.
}}

\VS{55}{\WJ{Când suflă un vânt de sud, voi spuneți: “Va fi o căldură toridă”, și așa se întâmplă.
}}

\VS{56}{\WJ{Voi, ipocriților! Voi știți să interpretați aspectul pământului și al cerului, dar cum se face că nu interpretați acest timp?
}}


\par }{\PP \VS{57}{\WJ{“De ce nu judecați voi înșivă ce este drept?
}}

\VS{58}{\WJ{Căci, atunci când mergi cu potrivnicul tău înaintea magistratului, încearcă cu sârguință pe drum să te eliberezi de el, ca nu cumva să nu te târască la judecător, iar judecătorul să te predea ofițerului și acesta să te arunce în închisoare.
}}

\VS{59}Îți
{\WJ{spun că nu vei ieși nicidecum de acolo până când nu vei plăti până la ultimul bănuț.”}}



\par }\Chap{13}{\PP \VerseOne{1}În același timp, erau de față unii care i-au spus despre galileenii al căror sânge Pilat îl amestecase cu jertfele lor.

\VS{2}Isus le-a răspuns:
{\WJ{“Credeți voi că acești galileeni erau mai păcătoși decât toți ceilalți galileeni, pentru că au suferit astfel de lucruri?
}}

\VS{3}{\WJ{Eu vă spun că nu, ci, dacă nu vă veți pocăi, veți pieri cu toții în același fel.
}}

\VS{4}{\WJ{Sau acei optsprezece peste care a căzut turnul din Siloam și i-a omorât — credeți că au fost niște păcătoși mai răi decât toți oamenii care locuiesc în Ierusalim?
}}

\VS{5}{\WJ{Vă spun că nu, dar, dacă nu vă veți pocăi, veți pieri cu toții în același fel.”
}}


\par }{\PP \VS{6}Și a spus pilda aceasta.
{\WJ{“Un om avea un smochin plantat în via sa; a venit să caute rod pe el și nu a găsit niciunul.
}}

\VS{7}{\WJ{El a zis viticultorului: “Iată, în acești trei ani am venit să caut rod pe acest smochin și n-am găsit niciunul. Tăiați-l! De ce risipește el pământul?”
}}

\VS{8}{\WJ{El a răspuns: “Doamne, lasă-l și anul acesta, până ce voi săpa în jurul lui și-l voi fertiliza.
}}

\VS{9}{\WJ{Dacă va da rod, bine; dacă nu, după aceea îl poți tăia'”.”
}}


\par }{\PP \VS{10}Învăța în una din sinagogi, în ziua de Sabat.

\VS{11}Și iată că era o femeie care avea un duh de slăbiciune de optsprezece ani. Era încovoiată și nu se putea în niciun fel îndrepta.

\VS{12}Când a văzut-o, Isus a chemat-o și i-a spus:
{\WJ{“Femeie, te-ai eliberat de boala ta”.
}}

\VS{13}Și-a pus mâinile peste ea, și imediat ea s-a ridicat în picioare și a slăvit pe Dumnezeu.


\par }{\PP \VS{14}Șeful sinagogii, indignat că Iisus vindecase în ziua Sabatului, a zis mulțimii: “Sunt șase zile în care trebuie să se lucreze. De aceea, veniți în acele zile și vindecați-vă, și nu în ziua de Sabat!”


\par }{\PP \VS{15}De aceea Domnul i-a răspuns: “{\WJ{Fățarnicilor! Oare nu-și eliberează fiecare dintre voi boul sau măgarul din grajd în ziua de Sabat și nu-l duce la apă?
}}

\VS{16}Oare
{\WJ{nu ar trebui ca această femeie, fiind o fiică a lui Avraam, pe care Satana a legat-o optsprezece ani lungi, să fie eliberată din această robie în ziua de Sabat?”
}}


\par }{\PP \VS{17}Pe când spunea el aceste lucruri, toți potrivnicii lui au fost dezamăgiți, și toată mulțimea s-a bucurat de toate lucrurile glorioase pe care le făcuse el.


\par }{\PP \VS{18}El a zis: “{\WJ{Cum este Împărăția lui Dumnezeu? Cu ce să o compar?
}}

\VS{19}Se
{\WJ{aseamănă cu un grăunte de muștar pe care l-a luat un om și l-a pus în grădina lui. Acesta a crescut și a devenit un copac mare, iar păsările cerului locuiesc în ramurile lui.”
}}


\par }{\PP \VS{20}Și a zis iarăși: “{\WJ{Cu ce să compar Împărăția lui Dumnezeu?
}}

\VS{21}{\WJ{Ea se aseamănă cu drojdia, pe care a luat-o o femeie și a ascuns-o în trei măsuri de făină, până ce a fost toată leșinată.”
}}


\par }{\PP \VS{22}Și mergea prin cetăți și sate, învățând și mergând până la Ierusalim.

\VS{23}Cineva i-a zis: “Doamne, sunt puțini cei care se mântuiesc?”

\par }{\PP El le-a zis:

\VS{24}“{\WJ{Străduiți-vă să intrați pe ușa cea strâmtă, căci vă spun că mulți vor căuta să intre și nu vor putea.
}}

\VS{25}{\WJ{După ce se va ridica stăpânul casei și va închide ușa, iar voi veți începe să stați afară și să bateți la ușă, zicând: “Doamne, Doamne, deschide-ne!”, atunci el vă va răspunde și vă va spune: “Nu vă cunosc și nu știu de unde veniți”.
}}

\VS{26}{\WJ{Atunci veți începe să spuneți: “Am mâncat și am băut în prezența ta și ai învățat pe străzile noastre.
}}

\VS{27}{\WJ{El va răspunde: 'Vă spun că nu știu de unde veniți. Depărtați-vă de la Mine, voi toți lucrătorii nelegiuirii'.
}}

\VS{28}{\WJ{Va fi plâns și scrâșnire de dinți când veți vedea pe Avraam, pe Isaac, pe Iacov și pe toți profeții în Împărăția lui Dumnezeu, iar voi înșivă veți fi aruncați afară.
}}

\VS{29}{\WJ{Ei vor veni de la răsărit, de la apus, de la nord și de la sud și se vor așeza în Împărăția lui Dumnezeu.
}}

\VS{30}{\WJ{Iată, unii dintre cei din urmă vor fi cei dintâi, iar unii dintre cei dintâi vor fi cei din urmă.”
}}


\par }{\PP \VS{31}În aceeași zi, au venit niște farisei și I-au zis: “Ieși de aici și pleacă, căci Irod vrea să Te omoare.”


\par }{\PP \VS{32}El le-a zis: “{\WJ{Duceți-vă și spuneți vulpii aceleia: “Iată că astăzi și mâine scot demoni și fac vindecări, iar a treia zi îmi termin misiunea.
}}

\VS{33}Cu toate
{\WJ{acestea, trebuie să-mi continui drumul astăzi, mâine și poimâine și poimâine, căci nu se poate ca un profet să piară în afara Ierusalimului'.
}}


\par }{\PP \VS{34}{\WJ{“Ierusalime, Ierusalime, Ierusalime, tu care ucizi pe prooroci și ucizi cu pietre pe cei trimiși la ea! De câte ori am vrut să-ți adun copiii, ca o găină care-și adună puii sub aripi, și tu ai refuzat!
}}

\VS{35}{\WJ{Iată, casa ta este lăsată pustie pentru tine. Vă spun că nu mă veți vedea până când nu veți spune: “Binecuvântat este cel care vine în numele Domnului!”.”
}}



\par }\Chap{14}{\PP \VerseOne{1}Când a intrat în casa unuia dintre fruntașii fariseilor, în ziua de Sabat, ca să mănânce pâine, ei Îl urmăreau.

\VS{2}Iată că un om care avea hidropizie era în fața lui.

\VS{3}Isus, răspunzând, a vorbit cu juriștii și cu fariseii și le-a zis: “{\WJ{Este îngăduit să vindeci în Sabat?”
}}


\par }{\PP \VS{4}Dar ei au tăcut.

\par }{\PP L-a luat, l-a vindecat și l-a lăsat să plece.

\VS{5}El le-a răspuns: “{\WJ{Cine dintre voi, dacă fiul său sau boul său ar cădea într-o fântână, nu l-ar scoate imediat într-o zi de Sabat?”
}}


\par }{\PP \VS{6}Nu puteau să-i răspundă la aceste lucruri.


\par }{\PP \VS{7}El a spus o pildă celor invitați, când a observat că ei își aleg locurile cele mai bune și le-a zis:

\VS{8}{\WJ{“Când sunteți invitați de cineva la o nuntă, nu vă așezați pe locul cel mai bun, căci poate că va fi invitat cineva mai de seamă decât voi,
}}

\VS{9}iar cel
{\WJ{care v-a invitat pe amândoi va veni și vă va spune: “Faceți loc pentru acesta”. Atunci ați începe, cu rușine, să ocupați locul cel mai de jos.
}}

\VS{10}{\WJ{Dar când ești invitat, du-te și stai în locul cel mai de jos, pentru ca, atunci când va veni cel care te-a invitat, să-ți spună: “Prietene, urcă-te mai sus”. Atunci vei fi onorat în prezența tuturor celor care stau la masă cu tine.
}}

\VS{11}{\WJ{Căci oricine se înalță va fi smerit, iar oricine se smerește va fi înălțat.”
}}


\par }{\PP \VS{12}Și a mai zis celui care-l invitase: “{\WJ{Când faci o cină sau un prânz, nu chema pe prietenii tăi, nici pe frații tăi, nici pe rudele tale, nici pe vecinii tăi bogați, căci poate că și ei îți vor întoarce favoarea și-ți vor plăti.
}}

\VS{13}Ci,
{\WJ{când faci un ospăț, invită-i pe săraci, pe mutilați, pe șchiopi sau pe orbi;
}}

\VS{14}și vei
{\WJ{fi binecuvântat, pentru că ei nu au resurse să-ți dea înapoi. Căci veți fi răsplătiți la învierea celor drepți.”
}}


\par }{\PP \VS{15}Unul din cei ce ședeau la masă cu El, auzind aceste lucruri, i-a zis: “Ferice de cel ce va petrece în Împărăția lui Dumnezeu!”


\par }{\PP \VS{16}Dar El i-a zis: “{\WJ{Un om a făcut o cină mare și a invitat multă lume.
}}

\VS{17}{\WJ{La ora cinei, a trimis pe slujitorul său să le spună celor invitați: “Veniți, căci totul este gata acum.
}}

\VS{18}{\WJ{Toți, ca unul singur, au început să se scuze.
}}

\par }{\PP {\WJ{“Cel dintâi i-a spus: “Am cumpărat un câmp și trebuie să mă duc să-l văd. Te rog să mă scuzi'.
}}


\par }{\PP \VS{19}{\WJ{Altul a zis: “Am cumpărat cinci juguri de boi și trebuie să mă duc să le încerc. Te rog să mă scuzi.''
}}


\par }{\PP \VS{20}{\WJ{Altul a zis: “M-am căsătorit cu o nevastă și nu pot să vin.
}}


\par }{\PP \VS{21}Și a
{\WJ{venit robul acela și a spus aceste lucruri domnului său. Atunci stăpânul casei, mâniindu-se, a zis robului său: “Ieși repede pe străzile și pe ulițele cetății și adu pe cei săraci, schilozi, orbi și șchiopi.
}}


\par }{\PP \VS{22}{\WJ{Robul a zis: “Doamne, s-a făcut cum ai poruncit și mai este loc.
}}


\par }{\PP \VS{23}{\WJ{Domnul a zis robului: “Ieși pe drumuri și pe ulițe și silește-i să intre, ca să se umple casa mea.
}}

\VS{24}{\WJ{Căci îți spun că niciunul dintre acei oameni care au fost invitați nu va gusta din cina mea.””
}}


\par }{\PP \VS{25}O mare mulțime de oameni mergeau cu El. El s-a întors și le-a zis:

\VS{26}{\WJ{“Dacă vine cineva la Mine și nu face abstracție de tatăl său, de mama sa, de soția sa, de copii, de frați și de surori, da, și chiar și de viața sa, nu poate fi ucenicul Meu.
}}

\VS{27}{\WJ{Oricine nu-și poartă propria cruce și nu vine după mine, nu poate fi discipolul meu.
}}

\VS{28}{\WJ{Căci care dintre voi, dorind să zidească un turn, nu stă mai întâi și numără costurile, ca să vadă dacă are destul pentru a-l termina?
}}

\VS{29}Sau
{\WJ{poate că, după ce a pus temelia și nu poate să o termine, toți cei care îl văd încep să râdă de el,
}}

\VS{30}{\WJ{spunând: “Omul acesta a început să zidească și nu a putut să termine”.
}}

\VS{31}{\WJ{Sau ce rege, când merge să se întâlnească cu un alt rege în război, nu se așează mai întâi și nu se gândește dacă este în stare cu zece mii de oameni să-l înfrunte pe cel care vine împotriva lui cu douăzeci de mii?
}}

\VS{32}{\WJ{Sau altfel, pe când celălalt este încă la mare distanță, trimite un trimis și cere condiții de pace.
}}

\VS{33}{\WJ{Așadar, oricine dintre voi care nu renunță la tot ceea ce are, nu poate fi discipolul meu.
}}


\par }{\PP \VS{34}{\WJ{“Sarea este bună, dar dacă sarea devine plată și fără gust, cu ce o condimentezi?
}}

\VS{35}{\WJ{Eanu este potrivită nici pentru pământ, nici pentru grămada de gunoi de grajd. Ea este aruncată la gunoi. Cine are urechi de auzit, să audă.”
}}



\par }\Chap{15}{\PP \VerseOne{1}Și toți vameșii și păcătoșii se apropiau de El ca să-L asculte.

\VS{2}Fariseii și cărturarii murmurau și ziceau: “Omul acesta îi primește pe păcătoși și mănâncă cu ei.”


\par }{\PP \VS{3}Și le-a spus pilda aceasta:

\VS{4}{\WJ{“Cine dintre voi, dacă ar avea o sută de oi și ar pierde una dintre ele, nu ar lăsa cele nouăzeci și nouă în pustiu și nu s-ar duce după cea pierdută, până o va găsi?
}}

\VS{5}După ce
{\WJ{o găsește, o poartă pe umeri, bucurându-se.
}}

\VS{6}{\WJ{Când se întoarce acasă, își cheamă prietenii și vecinii, spunându-le: “Bucurați-vă cu mine, căci am găsit oaia mea care era pierdută!”.
}}

\VS{7}{\WJ{Eu vă spun că și așa va fi mai multă bucurie în cer pentru un singur păcătos care se pocăiește, decât pentru nouăzeci și nouă de oameni drepți care nu au nevoie de pocăință.
}}


\par }{\PP \VS{8}{\WJ{“Sau ce femeie, dacă ar avea zece drahme, dacă ar pierde o drahmă, nu ar aprinde o lampă, nu ar mătura casa și nu ar căuta cu sârguință până ce ar găsi-o?
}}

\VS{9}După ce
{\WJ{o găsește, își cheamă prietenele și vecinele, spunând: “Bucurați-vă cu mine, căci am găsit drahma pe care o pierdusem!”.
}}

\VS{10}{\WJ{Tot așa, vă spun, există bucurie în prezența îngerilor lui Dumnezeu pentru un păcătos care se pocăiește.”
}}


\par }{\PP \VS{11}El a zis: “{\WJ{Un om avea doi fii.
}}

\VS{12}Cel
{\WJ{mai tânăr dintre ei a zis tatălui său: “Tată, dă-mi partea mea din averea ta. Așa că a împărțit între ei averea lui.
}}

\VS{13}{\WJ{Nu după multe zile, fiul cel mai tânăr a adunat toate acestea și a călătorit într-o țară îndepărtată. Acolo și-a risipit averea cu un trai desfrânat.
}}

\VS{14}După ce a
{\WJ{cheltuit totul, s-a ivit o foamete cruntă în țara aceea și a început să fie în nevoie.
}}

\VS{15}S-a
{\WJ{dus și s-a unit cu unul dintre cetățenii acelei țări, care l-a trimis pe câmpurile sale să hrănească porcii.
}}

\VS{16}{\WJ{El a vrut să-și umple burta cu păstăile pe care le mâncau porcii, dar nimeni nu i-a dat nimic.
}}

\VS{17}Când și-a
{\WJ{revenit, a zis: “Câți dintre angajații tatălui meu au pâine de prisos, iar eu mor de foame!
}}

\VS{18}Mă voi scula
{\WJ{și mă voi duce la tatăl meu și-i voi spune: “Tată, am păcătuit împotriva cerului și în fața ta.
}}

\VS{19}{\WJ{Nu mai sunt vrednic să fiu numit fiul tău. Fă-mă ca pe unul dintre robii tăi angajați!””'”.
}}


\par }{\PP \VS{20}{\WJ{“S-a sculat și a venit la tatăl său. Dar, pe când era încă departe, tatăl său l-a văzut și, mișcat de milă, a alergat, s-a aruncat la gâtul lui și l-a sărutat.
}}

\VS{21}{\WJ{Fiul i-a zis: “Tată, am păcătuit împotriva cerului și în fața ta. Nu mai sunt vrednic să mă numesc fiul tău'.
}}


\par }{\PP \VS{22}{\WJ{Dar tatăl a zis slugilor sale: “Scoateți cea mai bună haină și îmbrăcați-l. Puneți-i un inel la mână și sandale la picioare.
}}

\VS{23}{\WJ{Aduceți vițelul îngrășat, înjunghiați-l, să mâncăm și să sărbătorim;
}}

\VS{24}{\WJ{căci acesta, fiul meu, a fost mort și a înviat. Era pierdut și s-a regăsit”. Atunci au început să sărbătorească.
}}


\par }{\PP \VS{25}{\WJ{“Și fiul său cel mare era la câmp. Când s-a apropiat de casă, a auzit muzică și dansuri.
}}

\VS{26}{\WJ{A chemat la el pe unul dintre servitori și l-a întrebat ce se întâmplă.
}}

\VS{27}{\WJ{Acesta i-a spus: “A venit fratele tău, iar tatăl tău a omorât vițelul îngrășat, pentru că l-a primit înapoi sănătos și nevătămat.
}}

\VS{28}{\WJ{Dar el s-a mâniat și nu a vrut să intre. De aceea tatăl său a ieșit și l-a rugat.
}}

\VS{29}{\WJ{Dar el a răspuns tatălui său: “Iată, în acești mulți ani te-am slujit și niciodată nu am nesocotit o poruncă de-a ta, dar niciodată nu mi-ai dat un țap, ca să sărbătoresc cu prietenii mei.
}}

\VS{30}{\WJ{Dar când a venit acest fiu al tău, care ți-a devorat traiul cu prostituatele, ai ucis vițelul îngrășat pentru el.
}}


\par }{\PP \VS{31}{\WJ{El i-a zis: “Fiule, tu ești totdeauna cu mine și tot ce este al meu este al tău.
}}

\VS{32}{\WJ{Darse cuvenea să sărbătorim și să ne bucurăm, pentru că acesta, fratele tău, era mort și a înviat. Era pierdut și s-a găsit”.”
}}



\par }\Chap{16}{\PP \VerseOne{1}Și a mai zis ucenicilor Săi: “Era un
{\WJ{om bogat care avea un administrator. I s-a făcut o acuzație că acesta își risipește averile.
}}

\VS{2}{\WJ{El l-a chemat și i-a zis: “Ce este aceasta pe care o aud despre tine? Dă socoteală de administrarea ta, pentru că nu mai poți fi administrator'.
}}


\par }{\PP \VS{3}{\WJ{“Directorul și-a zis în sinea lui: “Ce voi face, dacă domnul meu îmi ia funcția de director? Nu am putere să sap. Mi-e rușine să cerșesc.
}}

\VS{4}{\WJ{Știu ce voi face, pentru ca, atunci când voi fi înlăturat de la conducere, să mă primească în casele lor.
}}

\VS{5}{\WJ{Chemând la el pe fiecare dintre datornicii stăpânului său, l-a întrebat pe primul: “Cât îi datorezi stăpânului meu?”
}}

\VS{6}{\WJ{El a răspuns: “O sută de batoane de ulei”. El i-a zis: 'Ia-ți nota de plată, așează-te repede și scrie cincizeci'.
}}

\VS{7}{\WJ{Apoi l-a întrebat pe altul: “Cât datorezi?”. El a răspuns: 'O sută de cors de grâu'. El i-a zis: “Ia-ți nota de plată și scrie optzeci”.
}}


\par }{\PP \VS{8}{\WJ{“Domnul său l-a lăudat pe administratorul necinstit pentru că a procedat cu înțelepciune, căci copiii acestei lumi sunt, în neamul lor, mai înțelepți decât copiii luminii.
}}

\VS{9}{\WJ{Eu vă spun: faceți-vă prieteni prin intermediul bogățiilor nedrepte, pentru ca, atunci când veți eșua, ei să vă primească în corturile veșnice.
}}

\VS{10}{\WJ{Cine este credincios în foarte puțin, este credincios și în mult. Cine este necinstit în foarte puțin, este necinstit și în mult.
}}

\VS{11}{\WJ{Așadar, dacă n-ați fost credincioși în bogăția nedreaptă, cine vă va încredința adevăratele bogății?
}}

\VS{12}{\WJ{Dacă n-ai fost credincios în ceea ce este al altuia, cine îți va da ceea ce este al tău?
}}

\VS{13}{\WJ{Nici un slujitor nu poate sluji la doi stăpâni, căci ori îl va urî pe unul și-l va iubi pe celălalt, ori se va ține de unul și-l va disprețui pe celălalt. Nu ești în stare să slujești lui Dumnezeu și lui Mamona.”
}}


\par }{\PP \VS{14}Fariseii, care erau iubitori de bani, au auzit și ei toate acestea și au luat în derâdere pe Isus.

\VS{15}El le-a zis: “{\WJ{Voi sunteți cei care vă justificați în fața oamenilor, dar Dumnezeu vă cunoaște inimile. Căci ceea ce este înălțat printre oameni este o urâciune în ochii lui Dumnezeu.
}}


\par }{\PP \VS{16}{\WJ{“Legea și proorocii au fost până la Ioan. De atunci se propovăduiește Vestea cea Bună a Împărăției lui Dumnezeu și fiecare își forțează intrarea în ea.
}}

\VS{17}{\WJ{Dar este mai ușor să treacă cerul și pământul, decât să cadă o mică lovitură de stilou din lege.
}}


\par }{\PP \VS{18}{\WJ{“Oricine divorțează de nevastă-sa și se căsătorește cu alta, săvârșește adulter. Cine se căsătorește cu una care a divorțat de un soț comite adulter.
}}


\par }{\PP \VS{19}“Era un
{\WJ{om bogat, îmbrăcat în purpură și în in subțire și trăind în fiecare zi în lux.
}}

\VS{20}{\WJ{Un cerșetor, numit Lazăr, a fost dus la poarta lui, plin de răni,
}}

\VS{21}{\WJ{și dorind să fie hrănit cu firimiturile care cădeau de la masa bogatului. Da, chiar și câinii veneau și îi lingeau rănile.
}}

\VS{22}{\WJ{Cerșetorul a murit și a fost dus de îngeri la sânul lui Avraam. Și bogatul a murit și el și a fost îngropat.
}}

\VS{23}{\WJ{În Hades, și-a ridicat ochii, fiind în chinuri, și a văzut departe pe Avraam și pe Lazăr la sânul lui.
}}

\VS{24}{\WJ{El a strigat și a zis: “Părinte Avraam, ai milă de mine și trimite-l pe Lazăr, ca să înmoaie vârful degetului în apă și să-mi răcorească limba! Căci sunt în chinuri în această flacără!”.
}}


\par }{\PP \VS{25}{\WJ{Dar Avraam a zis: “Fiule, adu-ți aminte că tu ai primit în viața ta cele bune, iar Lazăr, la fel, cele rele. Dar iată că el este acum mângâiat, iar tu ești în chinuri.
}}

\VS{26}În
{\WJ{afară de toate acestea, între noi și voi este fixată o prăpastie mare, încât cei care vor să treacă de aici la voi nu pot, și nimeni nu poate trece de acolo la noi'.
}}


\par }{\PP \VS{27}{\WJ{El a zis: “Te rog, părinte, să-l trimiți la casa tatălui meu,
}}

\VS{28}{\WJ{căci am cinci frați, ca să le dea mărturie, ca să nu ajungă și ei în locul acesta de chin.
}}


\par }{\PP \VS{29}{\WJ{Avraam i-a zis: “Ei au pe Moise și pe prooroci. Să-i asculte pe ei!”.
}}


\par }{\PP \VS{30}{\WJ{El a zis: “Nu, părinte Avraam, dar dacă se va duce la ei cineva din morți, se vor pocăi.
}}


\par }{\PP \VS{31}{\WJ{El i-a zis: “Dacă nu ascultă de Moise și de prooroci, nici dacă va învia unul din morți nu se vor lăsa convinși.”
}}



\par }\Chap{17}{\PP \VerseOne{1}El a zis ucenicilor: “{\WJ{Nu se poate să nu vină prilejuri de poticnire, ci vai de cel prin care vin!
}}

\VS{2}{\WJ{Ar fi mai bine pentru el să i se atârne de gât o piatră de moară și să fie aruncat în mare, decât să facă să se poticnească unul dintre acești micuți.
}}

\VS{3}{\WJ{Fiți atenți. Dacă fratele tău păcătuiește împotriva ta, mustră-l. Dacă se pocăiește, iartă-l.
}}

\VS{4}{\WJ{Dacă va păcătui împotriva ta de șapte ori într-o zi și de șapte ori se va întoarce, spunând: “Mă pocăiesc”, să-l ierți.”
}}


\par }{\PP \VS{5}Apostolii au zis Domnului: “Mărește-ne credința!”


\par }{\PP \VS{6}Domnul a zis: “{\WJ{Dacă ai avea credință ca un grăunte de muștar, ai spune acestui sicomor: “Dezrădăcinează-te și sădește-te în mare”, și el te-ar asculta.
}}

\VS{7}{\WJ{Dar cine dintre voi, având un rob care ară sau păzește oile, va spune, când vine de la câmp: 'Vino îndată și stai la masă'?
}}

\VS{8}{\WJ{Nu i-ar spune mai degrabă: 'Pregătește-mi cina, îmbracă-te cum se cuvine și servește-mă în timp ce mănânc și beau. După aceea vei mânca și vei bea'?”
}}

\VS{9}Îi
{\WJ{mulțumește oare acelui servitor pentru că a făcut ceea ce i s-a poruncit? Cred că nu.
}}

\VS{10}Tot
{\WJ{așa și voi, după ce ați făcut tot ce vi s-a poruncit, spuneți: 'Suntem niște servitori nevrednici. Ne-am făcut datoria'”.
}}


\par }{\PP \VS{11}Pe când se îndrepta spre Ierusalim, a trecut pe lângă hotarele Samariei și ale Galileii.

\VS{12}Pe când intra într-un sat, l-a întâmpinat zece bărbați leproși, care stăteau la distanță.

\VS{13}Ei și-au ridicat glasul și au zis: “Isuse, Învățătorule, ai milă de noi!”


\par }{\PP \VS{14}Când i-a văzut, le-a zis: “{\WJ{Duceți-vă și arătați-vă la preoți.” Pe măsură ce se duceau, au fost curățați.
}}

\VS{15}Unul dintre ei, când a văzut că s-a vindecat, s-a întors înapoi, slăvind pe Dumnezeu cu glas tare.

\VS{16}A căzut cu fața la picioarele lui Isus, mulțumindu-i; și era un samaritean.


\par }{\PP \VS{17}Isus a răspuns: “{\WJ{Nu au fost curățați cei zece? Dar unde sunt cei nouă?
}}

\VS{18}{\WJ{Nu s-a găsit niciunul care să se întoarcă să dea slavă lui Dumnezeu, în afară de acest străin?”
}}

\VS{19}Atunci i-a zis:
{\WJ{“Scoală-te și du-te. Credința ta te-a vindecat”.
}}


\par }{\PP \VS{20}Fiind întrebat de farisei când va veni Împărăția lui Dumnezeu, El le-a răspuns: “{\WJ{Împărăția lui Dumnezeu nu vine cu observație;
}}

\VS{21}{\WJ{nici nu vor spune: “Priviți, aici!” sau “Priviți, acolo!”, căci iată, Împărăția lui Dumnezeu este înăuntrul vostru”.
}}


\par }{\PP \VS{22}El a zis ucenicilor: “Vor
{\WJ{veni zile când veți dori să vedeți una din zilele Fiului Omului și nu o veți vedea.
}}

\VS{23}{\WJ{Ei vă vor spune: “Uite, aici!” sau “Uite, acolo!”. Nu plecați și nu mergeți după ei, nici nu vă luați după ei,
}}

\VS{24}{\WJ{pentru că, așa cum fulgerul, când strălucește dintr-o parte de sub cer, strălucește în altă parte de sub cer, așa va fi și Fiul Omului în ziua lui.
}}

\VS{25}{\WJ{Dar, mai întâi, trebuie să sufere multe și să fie respins de această generație.
}}

\VS{26}Cum a
{\WJ{fost în zilele lui Noe, așa va fi și în zilele Fiului Omului.
}}

\VS{27}{\WJ{Au mâncat, au băut, s-au căsătorit și s-au dat în căsătorie până în ziua în care Noe a intrat în corabie și a venit potopul și i-a distrus pe toți.
}}

\VS{28}Tot
{\WJ{așa, așa a fost și în zilele lui Lot: au mâncat, au băut, au cumpărat, au vândut, au plantat, au zidit;
}}

\VS{29}{\WJ{dar în ziua în care Lot a ieșit din Sodoma, a plouat foc și sulf din cer și i-a nimicit pe toți.
}}

\VS{30}La
{\WJ{fel se va întâmpla și în ziua în care se va arăta Fiul Omului.
}}

\VS{31}{\WJ{În ziua aceea, cel care va fi pe acoperișul casei și cu bunurile sale în casă, să nu se coboare să le ia. Cel care va fi pe câmp, la fel, să nu se întoarcă înapoi.
}}

\VS{32}Adu-ți
{\WJ{aminte de soția lui Lot!
}}

\VS{33}{\WJ{Cine caută să-și salveze viața o pierde, dar cine își pierde viața o păstrează.
}}

\VS{34}{\WJ{Vă spun că, în noaptea aceea, vor fi doi oameni în același pat. Unul va fi luat și celălalt va fi lăsat.
}}

\VS{35}{\WJ{Vor fi doi care vor măcina grâu împreună. Unul va fi luat și celălalt va rămâne.”
}}

\VS{36}\FTNT{*}{{\FR 17:36 }{\FT{Unele manuscrise grecești adaugă: „Doi vor fi pe câmp: unul luat, iar celălalt lăsat”.}}}


\par }{\PP \VS{37}Ei, răspunzând, L-au întrebat: “Unde, Doamne?”

\par }{\PP El le-a spus: “{\WJ{Unde este trupul, acolo se vor aduna și vulturii”.
}}



\par }\Chap{18}{\PP \VerseOne{1}Și le-a mai spus o pildă, ca să se roage mereu și să nu se dea bătut,

\VS{2}zicând: “{\WJ{Într-o cetate era un judecător care nu se temea de Dumnezeu și nu avea respect față de oameni.
}}

\VS{3}{\WJ{În acel oraș era o văduvă care venea adesea la el și îi spunea: “Apără-mă de adversarul meu!”.
}}

\VS{4}{\WJ{El nu a vrut pentru o vreme, dar după aceea și-a zis: “Deși nu mă tem de Dumnezeu și nici nu respect pe om,
}}

\VS{5}{\WJ{totuși, pentru că această văduvă mă deranjează, o voi apăra, altfel mă va obosi prin venirea ei continuă.””
}}


\par }{\PP \VS{6}Domnul a zis: “{\WJ{Ascultă ce spune judecătorul nedrept.
}}

\VS{7}{\WJ{Nu-i va răzbuna oare Dumnezeu pe aleșii Săi care strigă la El zi și noapte și totuși are răbdare cu ei?
}}

\VS{8}{\WJ{Eu vă spun că îi va răzbuna repede. Cu toate acestea, când va veni Fiul Omului, va găsi el credință pe pământ?”
}}


\par }{\PP \VS{9}Și a spus pilda aceasta unor oameni care erau convinși de dreptatea lor și care disprețuiau pe toți ceilalți:

\VS{10}{\WJ{“Doi oameni s-au urcat în templu să se roage: unul era fariseu, iar celălalt vameș.
}}

\VS{11}{\WJ{Fariseul stătea în picioare și se ruga de unul singur, așa: 'Doamne, îți mulțumesc că nu sunt ca ceilalți oameni: șantajiști, nedrepți, adulteri, sau chiar ca acest vameș.
}}

\VS{12}{\WJ{Eu postesc de două ori pe săptămână. Dau zeciuială din tot ce primesc'.
}}

\VS{13}{\WJ{Dar vameșul, care stătea departe, nici măcar nu voia să-și ridice ochii spre cer, ci își bătea pieptul, zicând: 'Doamne, ai milă de mine, păcătosul!
}}

\VS{14}{\WJ{Vă spun că acesta s-a coborât la casa lui mai degrabă îndreptățit decât celălalt; căci oricine se înalță va fi smerit, dar cel ce se smerește va fi înălțat.”
}}


\par }{\PP \VS{15}Îi aduceau și pruncii lor, ca să se atingă de ei. Dar, văzând aceasta, discipolii i-au mustrat.

\VS{16}Isus i-a chemat și le-a zis: “{\WJ{Lăsați copiii să vină la Mine și nu-i împiedicați, căci Împărăția lui Dumnezeu este a unora ca aceștia.
}}

\VS{17}{\WJ{Mai mult ca sigur, vă spun că oricine nu primește Împărăția lui Dumnezeu ca un copilaș, nu va intra nicidecum în ea.”
}}


\par }{\PP \VS{18}Un domnitor L-a întrebat: “Învățătorule, ce să fac ca să moștenesc viața veșnică?”


\par }{\PP \VS{19}Isus l-a întrebat: “{\WJ{De ce Mă numești bun? Nimeni nu este bun, în afară de unul singur: Dumnezeu.
}}

\VS{20}{\WJ{Tu cunoști poruncile: Să nu comiți adulter, să nu ucizi, să nu furi, să nu depui mărturie mincinoasă, să cinstești pe tatăl tău și pe mama ta.”
}}


\par }{\PP \VS{21}El a zis: “Am observat toate aceste lucruri din tinerețea mea.”


\par }{\PP \VS{22}Isus, auzind acestea, i-a zis: “{\WJ{Un singur lucru îți lipsește încă. Vinde tot ce ai și dă-le săracilor. Atunci vei avea o comoară în ceruri; apoi vino și urmează-mă.”
}}


\par }{\PP \VS{23}Dar când a auzit acestea, s-a întristat foarte tare, căci era foarte bogat.


\par }{\PP \VS{24}Isus, văzând că se întrista foarte tare, a zis: “{\WJ{Cât de greu intră în Împărăția lui Dumnezeu cei ce au bogății!
}}

\VS{25}{\WJ{Căci este mai ușor pentru o cămilă să intre prin urechea unui ac decât pentru un om bogat să intre în Împărăția lui Dumnezeu.”
}}


\par }{\PP \VS{26}Și cei ce au auzit au zis: “Atunci cine poate fi mântuit?”


\par }{\PP \VS{27}Dar El a zis: “Ceea ce
{\WJ{este cu neputință la oameni este cu putință la Dumnezeu.”
}}


\par }{\PP \VS{28}Petru a zis: “Uite, am lăsat totul și am venit după Tine.”


\par }{\PP \VS{29}El le-a zis: “{\WJ{Adevărat vă spun că nu este nimeni care a lăsat casă, sau nevastă, sau frați, sau părinți, sau copii, pentru Împărăția lui Dumnezeu,
}}

\VS{30}{\WJ{care să nu primească de multe ori mai mult în acest timp și în lumea viitoare viața veșnică.”
}}


\par }{\PP \VS{31}Apoi a luat pe cei doisprezece deoparte și le-a zis: “{\WJ{Iată, ne vom sui la Ierusalim și se vor împlini toate cele scrise prin prooroci despre Fiul Omului.
}}

\VS{32}{\WJ{Căci El va fi predat neamurilor, va fi batjocorit, tratat cu rușine și scuipat.
}}

\VS{33}{\WJ{Îl vor biciui și îl vor ucide. A treia zi, va învia.”
}}


\par }{\PP \VS{34}Ei n-au înțeles nimic din toate acestea. Acest cuvânt le-a fost ascuns și nu au înțeles ce se spunea.


\par }{\PP \VS{35}Pe când se apropia de Ierihon, un orb stătea pe drum și cerșea.

\VS{36}Auzind o mulțime care trecea, a întrebat ce înseamnă aceasta.

\VS{37}Ei i-au spus că Isus din Nazaret trecea pe acolo.

\VS{38}El a strigat: “Isuse, Fiul lui David, ai milă de mine!”

\VS{39}Cei care mergeau în frunte îl mustrau, ca să tacă; dar el striga și mai tare: “Fiul lui David, ai milă de mine!”


\par }{\PP \VS{40}Isus, stând în picioare, a poruncit să fie adus la El. După ce s-a apropiat, l-a întrebat:

\VS{41}{\WJ{“Ce vrei să fac?”.
}}

\par }{\PP El a spus: “Doamne, ca să văd din nou”.


\par }{\PP \VS{42}Isus i-a zis: “{\WJ{Primește-ți vederea. Credința ta te-a vindecat”.
}}


\par }{\PP \VS{43}Și îndată și-a recăpătat vederea și a mers după el, slăvind pe Dumnezeu. Tot poporul, când a văzut, a lăudat pe Dumnezeu.



\par }\Chap{19}{\PP \VerseOne{1}El a intrat și trecea prin Ierihon.

\VS{2}Era acolo un om cu numele Zaheu. Era un mare vameș și era bogat.

\VS{3}Încerca să vadă cine este Isus și nu putea din cauza mulțimii, pentru că era scund.

\VS{4}A alergat mai departe și s-a urcat într-un sicomor ca să îl vadă, pentru că urma să treacă pe acolo.

\VS{5}Când a ajuns în acel loc, Isus s-a uitat în sus, l-a văzut și i-a zis:
{\WJ{“Zaheu, grăbește-te și coboară-te, căci astăzi trebuie să rămân în casa ta.”
}}

\VS{6}El s-a grăbit, a coborât și l-a primit cu bucurie.

\VS{7}Când l-au văzut, toți au murmurat, zicând: “A intrat să se cazeze la un om păcătos”.


\par }{\PP \VS{8}Zacheu a stat în picioare și a zis Domnului: “Iată, Doamne, jumătate din averea mea o dau săracilor. Dacă am pretins ceva pe nedrept de la cineva, îi dau înapoi de patru ori mai mult.”


\par }{\PP \VS{9}Isus i-a zis: “{\WJ{Astăzi a venit mântuirea în casa aceasta, pentru că și el este fiu al lui Avraam.
}}

\VS{10}{\WJ{Căci Fiul Omului a venit să caute și să salveze ceea ce era pierdut.”
}}


\par }{\PP \VS{11}Auzind ei aceste lucruri, a continuat și a spus o pildă, pentru că era aproape de Ierusalim și ei credeau că Împărăția lui Dumnezeu se va arăta imediat.

\VS{12}El a spus deci: “Un
{\WJ{oarecare nobil s-a dus într-o țară îndepărtată ca să primească pentru sine o împărăție și să se întoarcă.
}}

\VS{13}{\WJ{El a chemat zece slujitori de-ai săi, le-a dat zece monede de mina și le-a spus: “Faceți afaceri până când voi veni eu”.
}}

\VS{14}{\WJ{Dar cetățenii lui l-au urât și au trimis un trimis după el, spunând: 'Nu vrem ca acest om să domnească peste noi'.
}}


\par }{\PP \VS{15}După ce s-a întors, după ce a
{\WJ{primit împărăția, a poruncit să fie chemați la el acești slujitori, cărora le dăduse banii, ca să știe ce au câștigat prin negoțul lor.
}}

\VS{16}{\WJ{Cel dintâi a venit înaintea lui, spunând: “Doamne, mina ta a făcut încă zece mine.
}}


\par }{\PP \VS{17}{\WJ{El i-a zis: “Bine ai făcut, slugă bună! Pentru că ai fost găsit credincios cu foarte puțin, vei avea autoritate peste zece cetăți'.
}}


\par }{\PP \VS{18}Al
{\WJ{doilea a venit și a zis: “Mina ta, Doamne, a făcut cinci mine.
}}


\par }{\PP \VS{19}{\WJ{Și i-a zis: “Și tu vei fi peste cinci cetăți.
}}


\par }{\PP \VS{20}Și a
{\WJ{venit un altul, zicând: “Doamne, iată mina ta, pe care o țineam în batistă,
}}

\VS{21}{\WJ{căci mă temeam de tine, pentru că ești un om exigent. Tu ridici ceea ce n-ai pus și culegi ceea ce n-ai semănat'.
}}


\par }{\PP \VS{22}{\WJ{El i-a zis: “Din gura ta te voi judeca, rob rău, din gura ta te voi judeca! Știai că sunt un om exigent, că iau ce nu am pus și culeg ce nu am semănat.
}}

\VS{23}{\WJ{Atunci de ce nu mi-ai depus banii la bancă și, la venirea mea, aș fi putut câștiga dobândă?”
}}

\VS{24}{\WJ{El a zis celor care stăteau de față: “Luați-i mina și dați-o celui care are zece mine.
}}


\par }{\PP \VS{25}{\WJ{“I-au zis: “Doamne, are zece mine!”.
}}

\VS{26}{\WJ{“Căci Eu vă spun că oricui are, i se va da mai mult; dar celui care nu are, i se va lua chiar și ceea ce are.
}}

\VS{27}{\WJ{Dar aduceți aici pe vrăjmașii mei care n-au vrut să domnesc peste ei și ucideți-i înaintea mea.””
}}

\VS{28}După ce a spus acestea, a plecat înainte, urcând la Ierusalim.


\par }{\PP \VS{29}Când S-a apropiat de Betfaghe și de Betania, pe muntele care se numește Măslinul, a trimis pe doi dintre ucenicii Săi,

\VS{30}zicând: “Mergeți
{\WJ{în satul de dincolo, în care, când veți intra, veți găsi legat un mânz pe care nimeni nu a șezut vreodată. Dezlegați-l și aduceți-l.
}}

\VS{31}{\WJ{Dacă vă întreabă cineva: “De ce îl dezlegi?”, spuneți-i: “Domnul are nevoie de el”.”
}}


\par }{\PP \VS{32}Cei trimiși s-au dus și au găsit lucrurile așa cum le spusese.

\VS{33}Pe când dezlegau mânzul, stăpânii lui i-au întrebat: “De ce dezlegați mânzul?”

\VS{34}Ei au răspuns: “Domnul are nevoie de el.”

\VS{35}Apoi l-au adus la Isus. Și-au aruncat mantiile pe mânz și l-au așezat pe Isus pe ele.

\VS{36}Pe când mergea, ei și-au întins mantiile pe drum.


\par }{\PP \VS{37}Pe când se apropia de El, la coborârea de pe Muntele Măslinilor, toată mulțimea ucenicilor a început să se bucure și să laude pe Dumnezeu cu glas tare pentru toate minunile pe care le văzuseră,

\VS{38}zicând: “Binecuvântat este Împăratul care vine în numele Domnului! Pace în ceruri și glorie în cele mai înalte!”.


\par }{\PP \VS{39}Unii dintre fariseii din mulțime I-au zis: “Învățătorule, mustră pe ucenicii Tăi!”


\par }{\PP \VS{40}El le-a răspuns: “{\WJ{Vă spun că, dacă acestea ar tăcea, pietrele ar striga.”
}}


\par }{\PP \VS{41}Când s-a apropiat, a văzut cetatea și a plâns din cauza ei,

\VS{42}zicând: “{\WJ{Dacă ai fi știut astăzi ce este al păcii tale! Dar acum, ele sunt ascunse de ochii tăi.
}}

\VS{43}{\WJ{Căci vor veni zile când vrăjmașii tăi vor ridica o baricadă împotriva ta, te vor înconjura, te vor strânge din toate părțile și
}}

\VS{44}{\WJ{te vor trânti la pământ pe tine și pe copiii tăi din tine. Nu vor lăsa în tine piatră peste piatră, pentru că nu ai știut timpul vizitei tale.”
}}


\par }{\PP \VS{45}El a intrat în Templu și a început să alunge pe cei ce cumpărau și vindeau în el,

\VS{46}zicându-le: “{\WJ{Este scris: “Casa Mea este o casă de rugăciune}}”,
{\WJ{iar voi ați făcut din ea o peșteră de tâlhari!”.
}}


\par }{\PP \VS{47}El învăța în fiecare zi în Templu, dar preoții cei mai de seamă, cărturarii și cei mai de seamă din popor căutau să-L nimicească.

\VS{48}Nu găseau ce să facă, căci tot poporul se agăța de fiecare cuvânt pe care-l spunea.



\par }\Chap{20}{\PP \VerseOne{1}Într-una din zilele acelea, pe când învăța pe norod în Templu și propovăduia vestea cea bună, au venit la El preoții și cărturarii și bătrânii.

\VS{2}Ei l-au întrebat: “Spune-ne: cu ce autoritate faci aceste lucruri? Sau cine îți dă această autoritate?”


\par }{\PP \VS{3}El le-a răspuns: “{\WJ{Și eu vă voi pune o întrebare. Spuneți-mi:
}}

\VS{4}{\WJ{Botezul lui Ioan a fost din cer sau de la oameni?”
}}


\par }{\PP \VS{5}Ei se gândeau între ei și ziceau: “Dacă vom zice: “Din cer”, va zice: “De ce nu l-ați crezut?”

\VS{6}Dar dacă vom zice: “De la oameni”, tot poporul ne va ucide cu pietre, pentru că sunt încredințați că Ioan a fost prooroc”.

\VS{7}Ei au răspuns că nu știau de unde este.


\par }{\PP \VS{8}Isus le-a zis: “{\WJ{Nici Eu nu vă voi spune cu ce putere fac aceste lucruri.”
}}


\par }{\PP \VS{9}Și a început să spună poporului pilda aceasta: “Un
{\WJ{om a plantat o vie, a dat-o în arendă unor agricultori și a plecat în altă țară pentru multă vreme.
}}

\VS{10}{\WJ{La vremea potrivită, a trimis un slujitor la fermieri pentru a-și lua partea lui din roadele viei. Dar fermierii l-au bătut și l-au trimis acasă cu mâna goală.
}}

\VS{11}{\WJ{A mai trimis încă un servitor, dar și pe acesta l-au bătut și l-au tratat cu rușine, și l-au trimis cu mâna goală.
}}

\VS{12}{\WJ{A mai trimis încă un al treilea, dar l-au rănit și pe el și l-au aruncat afară.
}}

\VS{13}{\WJ{Stăpânul viei a zis: “Ce să fac? Voi trimite pe fiul meu iubit. Poate că, văzându-l, îl vor respecta'.
}}


\par }{\PP \VS{14}{\WJ{Iar țăranii, văzându-l, s-au gândit între ei și au zis: “Acesta este moștenitorul. Haideți să-l omorâm, ca moștenirea să fie a noastră.'
}}

\VS{15}{\WJ{Atunci l-au aruncat afară din vie și l-au ucis. Ce le va face, așadar, stăpânul viei?
}}

\VS{16}{\WJ{El va veni și îi va nimici pe acești agricultori și va da via altora.”
}}

\par }{\PP Când au auzit asta, au spus: “Să nu se întâmple asta niciodată!”.


\par }{\PP \VS{17}Dar El, uitându-se la ei, a zis: “{\WJ{Ce este deci ceea ce este scris?
}}

\par }{\Q {\WJ{“Piatra pe care au respins-o zidarii
}}

\par }{\QB {\WJ{a fost făcută piatra de temelie?
}}


\par }{\Q \VS{18}{\WJ{Oricine va cădea pe piatra aceea va fi sfărâmat}},

\par }{\QB {\WJ{dar va face praf pe oricine va cădea pe el.”
}}


\par }{\PP \VS{19}Preoții cei mai de seamă și cărturarii au căutat să pună mâna pe El chiar în ceasul acela, dar se temeau de popor, căci știau că El spusese pilda aceasta împotriva lor.

\VS{20}L-au urmărit și au trimis spioni, care se prefăceau că sunt drepți, ca să-l prindă în capcană cu ceva din ceea ce ar fi spus, ca să-l predea puterii și autorității guvernatorului.

\VS{21}Ei l-au întrebat: “Învățătorule, știm că tu spui și înveți ceea ce este drept și nu ești părtinitor față de nimeni, ci înveți cu adevărat calea lui Dumnezeu.

\VS{22}Ne este permis să plătim sau nu impozite Cezarului?”


\par }{\PP \VS{23}Dar El, înțelegând viclenia lor, le-a zis:
{\WJ{“Pentru ce Mă puneți la încercare?
}}

\VS{24}{\WJ{Arătați-mi un denar. A cui sunt chipul și inscripția de pe el?”
}}

\par }{\PP Ei au răspuns: “A lui Cezar”.


\par }{\PP \VS{25}Și le-a zis: “{\WJ{Dați Cezarului ce este al Cezarului și lui Dumnezeu ce este al lui Dumnezeu.”
}}


\par }{\PP \VS{26}N-au putut să-l prindă în cuvintele lui în fața poporului. Ei s-au mirat de răspunsul lui și au tăcut.

\VS{27}Unii dintre saduchei au venit la el, cei care neagă că există o înviere.

\VS{28}L-au întrebat: “Învățătorule, Moise ne-a scris că, dacă fratele unui om moare având o soție, iar acesta nu are copii, fratele său trebuie să ia soția și să crească copii pentru fratele său.

\VS{29}Așadar, erau șapte frați. Primul și-a luat o soție și a murit fără copii.

\VS{30}Al doilea a luat-o de soție, și a murit fără copii.

\VS{31}Cel de-al treilea a luat-o și, la fel, toți cei șapte nu au lăsat copii și au murit.

\VS{32}După aceea a murit și femeia.

\VS{33}Prin urmare, la înviere, a cui soție dintre ei va fi ea? Căci cei șapte au avut-o ca soție.”


\par }{\PP \VS{34}Isus le-a zis: “{\WJ{Copiii veacului acestuia se căsătoresc și sunt dați în căsătorie.
}}

\VS{35}{\WJ{Dar cei care sunt considerați vrednici să ajungă la vârsta aceea și la învierea din morți nu se căsătoresc și nici nu sunt dați în căsătorie.
}}

\VS{36}{\WJ{Căci ei nu mai pot muri, pentru că sunt asemenea îngerilor și sunt copii ai lui Dumnezeu, fiind copii ai învierii.
}}

\VS{37}{\WJ{Dar că morții înviază, chiar Moise a arătat-o la tufiș, când L-a numit pe Domnul “Dumnezeul lui Avraam, Dumnezeul lui Isaac și Dumnezeul lui Iacov”.
}}

\VS{38}{\WJ{Or, El nu este Dumnezeul morților, ci al celor vii, căci pentru El toți sunt vii.”
}}


\par }{\PP \VS{39}Unii din cărturari au răspuns: “Învățătorule, bine vorbești.”

\VS{40}Nu au îndrăznit să-i mai pună întrebări.


\par }{\PP \VS{41}El le-a zis: “{\WJ{Pentru ce zic că Hristosul este fiul lui David?
}}

\VS{42}{\WJ{Însuși David spune în cartea Psalmilor
}}

\par }{\Q {\WJ{“Domnul a spus Domnului meu,
}}

\par }{\QB {\WJ{“Stai la dreapta mea,
}}


\par }{\QB \VS{43}{\WJ{până ce voi face din vrăjmașii tăi scăunelul picioarelor tale.”
}}


\par }{\PP \VS{44}{\WJ{“David îl numește deci Domn, și cum este el fiul său?”
}}


\par }{\PP \VS{45}În auzul întregului norod, a zis ucenicilor Săi:

\VS{46}“Feriți-vă
{\WJ{de cărturarii aceia cărora le place să umble în haine lungi, și care iubesc să fie salutați în piețe, să aibă cele mai bune locuri în sinagogi și cele mai bune locuri la ospețe;
}}

\VS{47}{\WJ{care mănâncă casele văduvelor și care, pentru a se preface, fac rugăciuni lungi. Aceștia vor primi o condamnare mai mare”.
}}



\par }\Chap{21}{\PP \VerseOne{1}Și, ridicându-și privirea, a văzut pe cei bogați care puneau darurile lor în vistierie.

\VS{2}A văzut o văduvă săracă care a aruncat două monede mici de aramă.

\VS{3}El a zis: “{\WJ{Adevărat vă spun că această văduvă săracă a pus mai mult decât toți aceștia,
}}

\VS{4}{\WJ{căci toți aceștia au pus daruri pentru Dumnezeu din belșugul lor, dar ea, din sărăcia ei, a pus tot ce avea ca să trăiască.”
}}


\par }{\PP \VS{5}Pe când vorbeau unii despre templu și despre cum era împodobit cu pietre frumoase și cu daruri, a zis:

\VS{6}“Cât
{\WJ{despre lucrurile acestea pe care le vedeți, vor veni zile în care nu va mai rămâne aici piatră pe piatră care să nu fie dărâmată.”
}}


\par }{\PP \VS{7}Ei L-au întrebat: “Învățătorule, când se vor întâmpla aceste lucruri? Care este semnul că aceste lucruri urmează să se întâmple?”


\par }{\PP \VS{8}Și a zis: “Luați
{\WJ{seama să nu vă lăsați rătăciți, căci mulți vor veni în numele Meu, zicând: “Eu sunt” și “Vremea este aproape”. De aceea, nu-i urmați.
}}

\VS{9}{\WJ{Când auziți de războaie și de tulburări, nu vă înspăimântați, pentru că aceste lucruri trebuie să se întâmple mai întâi, dar sfârșitul nu va veni imediat.”
}}


\par }{\PP \VS{10}Apoi le-a zis: “{\WJ{Neamul se va ridica împotriva neamului și împărăția împotriva împărăției.
}}

\VS{11}{\WJ{Vor fi cutremure mari, foamete și molime în diferite locuri. Vor fi spaime și semne mari din cer.
}}

\VS{12}{\WJ{Dar, înainte de toate acestea, vor pune mâna pe voi și vă vor persecuta, vă vor da în sinagogi și în închisori, vă vor aduce înaintea împăraților și a guvernatorilor, din pricina numelui Meu.
}}

\VS{13}Se va dovedi
{\WJ{ca o mărturie pentru voi.
}}

\VS{14}{\WJ{Așadar, așezați-vă în inimile voastre să nu meditați dinainte cum să răspundeți,
}}

\VS{15}{\WJ{pentru că eu vă voi da o gură și o înțelepciune pe care toți adversarii voștri nu vor putea să le înfrunte sau să le contrazică.
}}

\VS{16}{\WJ{Veți fi predați chiar de părinți, de frați, de rude și de prieteni. Aceștia vor face ca unii dintre voi să fie uciși.
}}

\VS{17}Veți
{\WJ{fi urâți de toți oamenii din cauza numelui meu.
}}

\VS{18}{\WJ{Și nici un fir de păr din capul vostru nu va pieri.
}}


\par }{\PP \VS{19}“{\WJ{Prin stăruința voastră veți câștiga viața voastră.
}}


\par }{\PP \VS{20}{\WJ{Dar când veți vedea Ierusalimul înconjurat de oștiri, să știți că s-a apropiat pustiirea lui.
}}

\VS{21}{\WJ{Atunci, cei care sunt în Iudeea să fugă în munți. Cei care sunt în mijlocul ei să plece. Cei care sunt în țară să nu intre în ea.
}}

\VS{22}{\WJ{Căci acestea sunt zile de răzbunare, ca să se împlinească tot ce este scris.
}}

\VS{23}{\WJ{Vai de cele însărcinate și de cele care alăptează prunci în acele zile! Căci va fi mare strâmtorare în țară și mânie pentru acest popor.
}}

\VS{24}{\WJ{Ei vor cădea sub ascuțișul sabiei și vor fi duși prizonieri în toate neamurile. Ierusalimul va fi călcat în picioare de neamuri până la împlinirea vremurilor neamurilor.
}}


\par }{\PP \VS{25}“{\WJ{Vor fi semne în soare, în lună și în stele, și pe pământ va fi neliniște printre popoare, în neliniște din cauza vuietului mării și al valurilor;
}}

\VS{26}{\WJ{oamenii vor fi înnebuniți de frică și de așteptare pentru lucrurile care vor veni peste lume, căci puterile cerurilor vor fi zguduite.
}}

\VS{27}{\WJ{Atunci vor vedea pe Fiul Omului venind pe un nor cu putere și cu o mare slavă.
}}

\VS{28}{\WJ{Dar când vor începe să se întâmple aceste lucruri, priviți în sus și ridicați-vă capetele, pentru că răscumpărarea voastră este aproape.”
}}


\par }{\PP \VS{29}Și le-a spus o pildă.
{\WJ{“Priviți smochinul și toți copacii.
}}

\VS{30}{\WJ{Când sunt deja înmuguriți, voi vedeți și știți de la sine că vara este deja aproape.
}}

\VS{31}Tot
{\WJ{așa și voi, când vedeți că se întâmplă aceste lucruri, știți că Împărăția lui Dumnezeu este aproape.
}}

\VS{32}{\WJ{Mai mult ca sigur vă spun că generația aceasta nu va trece până când nu se vor împlini toate lucrurile.
}}

\VS{33}{\WJ{Cerul și pământul vor trece, dar cuvintele mele nu vor trece nicidecum.
}}


\par }{\PP \VS{34}{\WJ{Deci, luați seama, căci inimile voastre se vor împovăra de desfătări, de beții și de grijile vieții acesteia, și ziua aceea va veni pe neașteptate asupra voastră.
}}

\VS{35}{\WJ{Căci ea va veni ca un laț peste toți cei care locuiesc pe suprafața întregului pământ.
}}

\VS{36}{\WJ{De aceea, vegheați tot timpul, rugându-vă ca să fiți socotiți vrednici să scăpați de toate aceste lucruri care se vor întâmpla și să stați în picioare înaintea Fiului Omului.”
}}


\par }{\PP \VS{37}În fiecare zi, Isus învăța în Templu, iar noaptea ieșea să înnopteze pe muntele care se numește Măslinii.

\VS{38}Tot poporul venea dis-de-dimineață la el în templu ca să-l asculte.



\par }\Chap{22}{\PP \VerseOne{1}Se apropia sărbătoarea azimilor, care se cheamă Paștele.

\VS{2}Preoții cei mai de seamă și cărturarii căutau cum să-L omoare, pentru că se temeau de popor.


\par }{\PP \VS{3}Satana a intrat în Iuda, numit și Iscarioteanul, care era numărat între cei doisprezece.

\VS{4}El s-a dus și a vorbit cu preoții cei mai de seamă și cu căpeteniile despre cum să li-l predea.

\VS{5}Ei s-au bucurat și au fost de acord să-i dea bani.

\VS{6}El a consimțit și a căutat un prilej să li-l predea în absența mulțimii.


\par }{\PP \VS{7}A sosit ziua azimilor, în care trebuia să se jertfească Paștele.

\VS{8}Isus i-a trimis pe Petru și pe Ioan, zicând:
{\WJ{“Duceți-vă și pregătiți-ne Paștele, ca să mâncăm.”
}}


\par }{\PP \VS{9}Ei I-au zis: “Unde vrei să ne pregătim?”


\par }{\PP \VS{10}El le-a zis: “{\WJ{Iată, când veți intra în cetate, vă va întâmpina un om care duce un ulcior cu apă. Urmați-l în casa în care va intra el.
}}

\VS{11}{\WJ{Spuneți-i stăpânului casei: “Învățătorul vă spune: “Unde este camera de oaspeți, unde pot mânca Paștele cu discipolii mei?””.
}}

\VS{12}{\WJ{Îți va arăta o cameră mare, mobilată, în partea de sus. Fă pregătirile acolo”.
}}


\par }{\PP \VS{13}S-au dus, au găsit lucrurile așa cum le spusese Isus și au pregătit Paștele.


\par }{\PP \VS{14}Când a sosit ceasul, a șezut la masă cu cei doisprezece apostoli.

\VS{15}Și le-a zis:
{\WJ{“Am dorit foarte mult să mănânc acest Paște cu voi înainte de a suferi,
}}

\VS{16}{\WJ{căci vă spun că nu voi mai mânca în niciun fel din el până când nu se va împlini în Împărăția lui Dumnezeu.”
}}

\VS{17}A primit un pahar și, după ce a mulțumit, a zis:
{\WJ{“Luați și împărțiți-l între voi,
}}

\VS{18}{\WJ{pentru că vă spun că nu voi mai bea deloc din rodul viței de vie, până când nu va veni Împărăția lui Dumnezeu.”
}}


\par }{\PP \VS{19}A luat o pâine și, după ce a mulțumit, a frânt-o și le-a dat-o, zicând: “{\WJ{Acesta este trupul Meu, care se dă pentru voi. Faceți aceasta în amintirea mea”.
}}

\VS{20}La fel, după cină, a luat paharul, zicând:
{\WJ{“Paharul acesta este noul legământ în sângele Meu, care se varsă pentru voi.}}

\VS{21}{\WJ{Dar iată că mâna celui care mă trădează este cu mine pe masă.
}}

\VS{22}{\WJ{Într-adevăr, Fiul Omului merge așa cum a fost hotărât, dar vai de omul prin care este trădat!”
}}


\par }{\PP \VS{23}Și au început să se întrebe între ei care dintre ei era cel ce voia să facă lucrul acesta.


\par }{\PP \VS{24}Între ei s-a iscat și o ceartă, care dintre ei era considerat cel mai mare.

\VS{25}El le-a zis: “{\WJ{Împărații neamurilor îi stăpânesc, iar cei care au autoritate asupra lor sunt numiți “binefăcători”.
}}

\VS{26}{\WJ{Dar nu este așa cu voi. Mai degrabă, cel care este mai mare dintre voi să devină ca cel mai mic, iar cel care conduce, ca unul care slujește.
}}

\VS{27}{\WJ{Căci cine este mai mare, cel care stă la masă sau cel care slujește? Nu este oare cel care stă la masă? Dar eu sunt în mijlocul vostru ca unul care slujește.
}}


\par }{\PP \VS{28}{\WJ{Dar voi sunteți aceia care ați rămas cu mine în încercările mele.
}}

\VS{29}{\WJ{Eu vă dau vouă o împărăție, așa cum Tatăl Meu mi-a dat-o Mie,
}}

\VS{30}{\WJ{ca să mâncați și să beți la masa Mea în Împărăția Mea. Voi veți sta pe tronuri, judecând cele douăsprezece triburi ale lui Israel.”
}}


\par }{\PP \VS{31}Domnul a zis:
{\WJ{“Simon, Simon, iată că Satana a cerut să vă ia pe toți, ca să vă cerne ca pe grâu;
}}

\VS{32}{\WJ{dar Eu m-am rugat pentru tine, ca să nu-ți piardă credința. Tu, după ce te vei fi întors din nou, stabilește-i pe frații tăi”.
}}


\par }{\PP \VS{33}El I-a zis: “Doamne, sunt gata să merg cu Tine și la temniță și la moarte!”


\par }{\PP \VS{34}El a zis:
{\WJ{“Îți spun, Petru, că astăzi nu va cânta cocoșul până nu vei nega de trei ori că Mă cunoști.”
}}


\par }{\PP \VS{35}El le-a zis: “{\WJ{Când v-am trimis fără pungă, fără sac și fără sandale, v-a lipsit ceva?”
}}

\par }{\PP Ei au spus: “Nimic”.


\par }{\PP \VS{36}Apoi le-a zis: “{\WJ{Dar acum, cine are o pungă, s-o ia și un sac. Cine n-are, să-și vândă haina și să-și cumpere o sabie.
}}

\VS{37}{\WJ{Căci vă spun că trebuie să se împlinească încă în mine ceea ce este scris: “A fost socotit împreună cu călcătorii de lege”. Căci ceea ce mă privește pe mine se împlinește.”
}}


\par }{\PP \VS{38}Ei au zis: “Doamne, iată două săbii.”

\par }{\PP El le-a spus:
{\WJ{“Ajunge!”.
}}


\par }{\PP \VS{39}El a ieșit și s-a dus, după obiceiul său, pe Muntele Măslinilor. L-au urmat și discipolii Lui.

\VS{40}Când a ajuns la locul acela, le-a zis:
{\WJ{“Rugați-vă să nu intrați în ispită.”
}}


\par }{\PP \VS{41}S-a depărtat de ei ca la o aruncătură de piatră, a îngenuncheat și se ruga, zicând:

\VS{42}“{\WJ{Tată, dacă voiești, depărtează de la Mine paharul acesta. Cu toate acestea, nu voia mea, ci a Ta să se facă”.”
}}


\par }{\PP \VS{43}Un înger din cer i s-a arătat și l-a întărit.

\VS{44}Fiind în agonie, el se ruga cu mai multă stăruință. Sudoarea lui a devenit ca niște picături mari de sânge care cădeau pe pământ.


\par }{\PP \VS{45}După ce S-a sculat din rugăciunea Sa, a venit la ucenici și, găsindu-i adormiți din pricina durerii,

\VS{46}le-a zis:
{\WJ{“De ce dormiți? Sculați-vă și rugați-vă ca să nu intrați în ispită”.
}}


\par }{\PP \VS{47}Pe când vorbea el încă, a apărut o mulțime. Cel care se numea Iuda, unul dintre cei doisprezece, îi conducea. S-a apropiat de Isus ca să-l sărute.

\VS{48}Dar Isus i-a zis:
{\WJ{“Iuda, tu Îl trădezi pe Fiul Omului cu o sărutare?”
}}


\par }{\PP \VS{49}Cei ce erau în jurul lui, văzând ce avea să se întâmple, I-au zis: “Doamne, să batem cu sabia?”

\VS{50}Unul dintre ei a lovit pe slujitorul marelui preot și i-a tăiat urechea dreaptă.


\par }{\PP \VS{51}Dar Isus a răspuns: “{\WJ{Lasă-mă măcar să fac lucrul acesta, și, atingându-i urechea, l-a vindecat.
}}

\VS{52}Isus a zis preoților cei mai de seamă, căpeteniilor templului și bătrânilor, care veniseră împotriva Lui: “Ați
{\WJ{ieșit ca împotriva unui tâlhar, cu săbii și cu ciomege?
}}

\VS{53}{\WJ{Când eram zilnic cu voi în templu, nu ați întins mâinile împotriva mea. Dar acesta este ceasul vostru și puterea întunericului”.
}}


\par }{\PP \VS{54}L-au prins și l-au dus și l-au adus în casa marelui preot. Dar Petru îl urmărea de la distanță.

\VS{55}După ce au aprins un foc în mijlocul curții și au șezut împreună, Petru a șezut în mijlocul lor.

\VS{56}O slujnică l-a văzut cum ședea la lumină și, uitându-se cu atenție la el, a zis: “Și acesta era cu el.”


\par }{\PP \VS{57}El s-a lepădat de Isus, zicând: “Femeie, nu-L cunosc.”


\par }{\PP \VS{58}După puțină vreme, l-a văzut altcineva și i-a zis: “Și tu ești unul dintre ei!”

\par }{\PP Dar Petru a răspuns: “Omule, eu nu sunt!”


\par }{\PP \VS{59}După ce a trecut aproape un ceas, un altul a spus cu încredere: “Cu adevărat și acesta era cu El, căci este galileean.”


\par }{\PP \VS{60}Dar Petru a zis: “Omule, nu știu despre ce vorbești!” Imediat, în timp ce el încă vorbea, a cântat un cocoș.

\VS{61}Domnul s-a întors și s-a uitat la Petru. Atunci Petru și-a adus aminte de cuvântul Domnului, cum că i-a spus:
{\WJ{“Înainte să cânte cocoșul, te vei lepăda de mine de trei ori”.
}}

\VS{62}A ieșit afară și a plâns cu amar.


\par }{\PP \VS{63}Oamenii care îl țineau pe Isus își băteau joc de el și îl băteau.

\VS{64}După ce l-au legat la ochi, l-au lovit peste față și l-au întrebat: “Proorocește! Cine este cel care te-a lovit?”.

\VS{65}Și au mai spus multe alte lucruri împotriva lui, insultându-l.


\par }{\PP \VS{66}Când s-a făcut ziuă, s-a adunat adunarea bătrânilor poporului, preoții cei mai de seamă și cărturarii, și L-au dus pe Isus în sfatul lor, zicând:

\VS{67}“Dacă Tu ești Hristosul, spune-ne.”

\par }{\PP Dar El le-a zis: “Dacă
{\WJ{vă voi spune, nu veți crede;
}}

\VS{68}și dacă vă voi
{\WJ{întreba, nu-mi veți răspunde și nu mă veți lăsa să plec.
}}

\VS{69}De
{\WJ{acum încolo, Fiul Omului va ședea la dreapta puterii lui Dumnezeu”.
}}


\par }{\PP \VS{70}Și toți au zis: “Deci Tu ești Fiul lui Dumnezeu?”

\par }{\PP El le-a spus:
{\WJ{“Spuneți asta, pentru că așa sunt”.
}}


\par }{\PP \VS{71}Ei au zis: “Pentru ce mai avem nevoie de un martor? Pentru că noi înșine am auzit din gura Lui!”



\par }\Chap{23}{\PP \VerseOne{1}Toată ceata lor s-a sculat și L-au dus înaintea lui Pilat.

\VS{2}Ei au început să-l acuze, zicând: “Am găsit pe omul acesta pervertindu-și națiunea, interzicând plata impozitelor către Cezar și spunând că el însuși este Hristos, un rege.”


\par }{\PP \VS{3}Pilat L-a întrebat: “Tu ești Împăratul Iudeilor?”

\par }{\PP El i-a răspuns:
{\WJ{“Așa spui tu”.
}}


\par }{\PP \VS{4}Pilat a zis preoților de seamă și mulțimii: “Nu găsesc nici un temei de acuzare împotriva omului acesta.”


\par }{\PP \VS{5}Dar ei insistau, zicând: “El stârnește poporul și învață în toată Iudeea, începând din Galileea și până aici.”


\par }{\PP \VS{6}Pilat, când a auzit vorbindu-se de Galileea, a întrebat dacă omul acela este galileean.

\VS{7}Când a aflat că se afla sub jurisdicția lui Irod, l-a trimis la Irod, care se afla și el la Ierusalim în acele zile.


\par }{\PP \VS{8}Când a văzut Irod pe Isus, s-a bucurat foarte mult, căci de multă vreme voia să-L vadă, pentru că auzise multe lucruri despre El. Spera să vadă vreo minune făcută de el.

\VS{9}L-a întrebat cu multe cuvinte, dar el nu i-a dat niciun răspuns.

\VS{10}Preoții cei mai de seamă și cărturarii stăteau în picioare, acuzându-l cu vehemență.

\VS{11}Irod cu soldații săi îl umileau și își băteau joc de el. Îmbrăcându-l în haine luxoase, l-au trimis înapoi la Pilat.

\VS{12}Irod și Pilat s-au împrietenit chiar în ziua aceea, căci înainte de aceasta erau dușmani unul cu altul.


\par }{\PP \VS{13}Pilat a chemat pe preoții cei mai de seamă, pe dregători și pe popor,

\VS{14}și le-a zis: “Mi l-ați adus pe omul acesta ca pe unul care strică poporul și, după ce l-am cercetat înaintea voastră, n-am găsit nici un temei de acuzație împotriva omului acesta pentru cele de care îl acuzați.

\VS{15}Nici Irod, căci v-am trimis la el și vedeți, nu a făcut nimic vrednic de moarte.

\VS{16}De aceea îl voi pedepsi și îl voi elibera.”


\par }{\PP \VS{17}Și a trebuit să le dea drumul la un prizonier la sărbătoare.

\VS{18}Dar toți au strigat împreună și au zis: “Departe de omul acesta! Eliberați-ne pe Baraba!” —

\VS{19}unul care fusese aruncat în închisoare pentru o anumită revoltă în cetate și pentru crimă.


\par }{\PP \VS{20}Atunci Pilat le-a vorbit iarăși și a vrut să dea drumul lui Isus;

\VS{21}dar ei strigau și ziceau: “Răstignește! Răstignește-l!”


\par }{\PP \VS{22}A treia oară le-a zis: “De ce? Ce rău a făcut omul acesta? Nu am găsit în el nicio crimă capitală. De aceea îl voi pedepsi și îl voi elibera.”

\VS{23}Dar ei insistau cu glasuri puternice, cerând să fie răstignit. Vocile lor și ale preoților de seamă au învins.

\VS{24}Pilat a decretat să se facă ceea ce cereau ei.

\VS{25}L-a eliberat pe cel care fusese aruncat în închisoare pentru insurecție și crimă, pentru care cereau ei, dar pe Isus l-a predat voinței lor.


\par }{\PP \VS{26}După ce L-au dus, au apucat pe Simon din Cirene, care venea de la țară, și au pus crucea pe el, ca s-o ducă după Isus.

\VS{27}O mare mulțime de oameni îl urmau, inclusiv femei care și ele îl plângeau și îl jeleau.

\VS{28}Dar Isus, întorcându-se spre ele, le-a zis: “{\WJ{Fiice ale Ierusalimului, nu plângeți pentru mine, ci plângeți pentru voi și pentru copiii voștri.
}}

\VS{29}{\WJ{Căci iată că vin zile în care se va spune: “Ferice de cele sterpe, de pântecele care n-au născut niciodată și de sânii care n-au alăptat niciodată!”.
}}

\VS{30}{\WJ{Atunci vor începe să spună munților: “Cădeți peste noi!” și vor spune dealurilor: “Acoperiți-ne!”.
}}

\VS{31}{\WJ{Căci, dacă fac aceste lucruri în copacul verde, ce se va face în cel uscat?”
}}


\par }{\PP \VS{32}Au mai fost și alți doi criminali, care au fost duși împreună cu el la moarte.

\VS{33}Când au ajuns la locul numit “Craniul”, L-au răstignit acolo împreună cu criminalii, unul la dreapta și altul la stânga.


\par }{\PP \VS{34}Isus a zis:
{\WJ{“Tată, iartă-i, căci nu știu ce fac.”
}}

\par }{\PP Împărțind hainele lui între ei, au tras la sorți.

\VS{35}Poporul stătea și privea. Și conducătorii care erau cu ei își băteau joc de el, zicând: “A salvat pe alții. Să se salveze pe sine însuși, dacă acesta este Hristosul lui Dumnezeu, alesul lui!”


\par }{\PP \VS{36}Soldații își băteau joc de El, venind la El și oferindu-i oțet,

\VS{37}și zicând: “Dacă ești Împăratul Iudeilor, mântuiește-te!”


\par }{\PP \VS{38}De asemenea, deasupra lui era scrisă o inscripție cu litere grecești, latinești și ebraice: “Acesta este REGELE EVREILOR”.


\par }{\PP \VS{39}Unul dintre cei spânzurați L-a insultat, zicând: “Dacă ești Hristosul, mântuiește-te pe Tine și pe noi!”


\par }{\PP \VS{40}Dar celălalt a răspuns și, mustrându-l, i-a zis: “Nu te temi tu de Dumnezeu, când ești sub aceeași osândă?

\VS{41}Și noi, într-adevăr, suntem drepți, căci primim răsplata cuvenită pentru faptele noastre, dar omul acesta nu a făcut nimic rău.”

\VS{42}El i-a zis lui Isus: “Doamne, adu-ți aminte de mine când vei veni în Împărăția ta”.


\par }{\PP \VS{43}Isus i-a zis: “{\WJ{Adevărat îți spun că astăzi vei fi cu Mine în Rai.”
}}


\par }{\PP \VS{44}Era pe la ceasul al șaselea și s-a făcut întuneric peste toată țara până la ceasul al nouălea.

\VS{45}Soarele s-a întunecat, iar perdeaua templului s-a rupt în două.

\VS{46}Isus, strigând cu glas tare, a zis:
{\WJ{“Tată, în mâinile Tale Îmi dau duhul Meu!” După ce a spus aceasta, și-a dat ultima suflare.
}}


\par }{\PP \VS{47}Când a văzut centurionul ce se făcuse, a slăvit pe Dumnezeu, zicând: “Cu adevărat, acesta era un om drept.”

\VS{48}Toate mulțimile care se adunaseră să vadă acest lucru, când au văzut ce s-a făcut, s-au întors acasă bătându-se pe piept.

\VS{49}Toți cunoscuții lui și femeile care îl însoțeau din Galileea stăteau deoparte, privind aceste lucruri.


\par }{\PP \VS{50}Și iată că era un om numit Iosif, care făcea parte din consiliu, un om bun și drept

\VS{51}(nu consimțise la sfatul și la fapta lor), din Arimateea, cetate a iudeilor, care aștepta și el Împărăția lui Dumnezeu.

\VS{52}Acest om s-a dus la Pilat și a cerut trupul lui Isus.

\VS{53}Acesta l-a dat jos, l-a înfășurat într-o pânză de in și l-a pus într-un mormânt tăiat în piatră, unde nu fusese pus nimeni niciodată.

\VS{54}Era ziua Pregătirii și se apropia Sabatul.

\VS{55}Femeile care veniseră cu el din Galileea au venit după el și au văzut mormântul și cum fusese pus trupul Lui.

\VS{56}Ele s-au întors și au pregătit miresme și miruri. În Sabat s-au odihnit conform poruncii.



\par }\Chap{24}{\PP \VerseOne{1}Dar în prima zi a săptămânii, în zorii zilei, au venit la mormânt, împreună cu alții, aducând miresmele pe care le pregătiseră.

\VS{2}Au găsit piatra rostogolită de pe mormânt.

\VS{3}Au intrat înăuntru și n-au găsit trupul Domnului Isus.

\VS{4}În timp ce erau foarte nedumeriți de acest lucru, iată că au apărut lângă ei doi bărbați îmbrăcați în haine orbitoare.

\VS{5}Îngrozindu-se, s-au plecat cu fața la pământ.

\par }{\PP Oamenii le-au zis: “De ce căutați pe cei vii printre cei morți?

\VS{6}El nu este aici, ci a înviat. Vă amintiți ce v-a spus când era încă în Galileea,

\VS{7}spunând că Fiul Omului trebuie să fie dat în mâinile oamenilor păcătoși, să fie răstignit și a treia zi să învieze?”


\par }{\PP \VS{8}Ei și-au adus aminte de cuvintele Lui,

\VS{9}s-au întors de la mormânt și au spus toate acestea celor unsprezece și tuturor celorlalți.

\VS{10}Ele erau Maria Magdalena, Ioana și Maria, mama lui Iacov. Celelalte femei care erau cu ele au povestit aceste lucruri apostolilor.

\VS{11}Aceste cuvinte li s-au părut a fi un nonsens și nu le-au crezut.

\VS{12}Dar Petru s-a ridicat și a alergat la mormânt. Aplecându-se și uitându-se înăuntru, a văzut fâșiile de pânză așezate singure și a plecat acasă, întrebându-se ce se întâmplase.


\par }{\PP \VS{13}Și iată că doi dintre ei mergeau chiar în ziua aceea într-un sat numit Emaus, care era la șaizeci de stadii de Ierusalim.

\VS{14}Și vorbeau unul cu altul despre toate aceste lucruri care se întâmplaseră.

\VS{15}În timp ce vorbeau și se întrebau împreună, Isus însuși s-a apropiat și a mers cu ei.

\VS{16}Dar ochii lor erau împiedicați să-L recunoască.

\VS{17}El le-a zis:
{\WJ{“Despre ce vorbiți în timp ce mergeți și sunteți triști?”
}}


\par }{\PP \VS{18}Unul dintre ei, numit Cleopa, i-a răspuns: “Ești tu singurul străin din Ierusalim care nu știe ce s-a întâmplat acolo în zilele acestea?”


\par }{\PP \VS{19}El le-a zis:
{\WJ{“Ce anume?”
}}

\par }{\PP Ei I-au zis: “Cele despre Iisus Nazarineanul, care a fost un prooroc puternic cu fapta și cu cuvântul înaintea lui Dumnezeu și a întregului popor,

\VS{20}și despre cum preoții cei mai de seamă și căpeteniile noastre L-au dat spre osândă la moarte și L-au răstignit.

\VS{21}Dar noi speram că el este cel care va răscumpăra pe Israel. Da, și, pe lângă toate acestea, este acum a treia zi de când s-au întâmplat aceste lucruri.

\VS{22}De asemenea, ne-au uimit niște femei din anturajul nostru, care au ajuns mai devreme la mormânt;

\VS{23}și, când nu i-au găsit trupul, au venit spunând că au văzut și ele o viziune a unor îngeri, care spuneau că este viu.

\VS{24}Unii dintre noi s-au dus la mormânt și l-au găsit așa cum au spus femeile, dar nu l-au văzut.”


\par }{\PP \VS{25}El le-a zis: “{\WJ{Popor nebun și încet la inimă ca să creadă în tot ce au spus proorocii!
}}

\VS{26}{\WJ{Nu trebuia oare Hristosul să sufere aceste lucruri și să intre în gloria Sa?”
}}

\VS{27}Pornind de la Moise și de la toți profeții, le-a explicat în toate Scripturile ceea ce se referă la el însuși.


\par }{\PP \VS{28}Au ajuns aproape de satul în care se duceau, iar el s-a făcut că vrea să meargă mai departe.


\par }{\PP \VS{29}Ei l-au îndemnat, zicând: “Rămâi cu noi, căci este aproape seară și ziua este aproape de sfârșit.”

\par }{\PP A intrat să stea cu ei.

\VS{30}După ce s-a așezat la masă cu ei, a luat pâinea și a mulțumit. A frânt-o și le-a dat-o.

\VS{31}Ochii lor s-au deschis și l-au recunoscut; apoi a dispărut din ochii lor.

\VS{32}Ei și-au zis unul altuia: “Nu ne ardea inima în noi când ne vorbea pe drum și când ne deschidea Scripturile?”

\VS{33}Ei s-au sculat chiar în acel ceas, s-au întors la Ierusalim și i-au găsit pe cei unsprezece adunați împreună cu cei care erau cu ei,

\VS{34}spunând: “Domnul a înviat cu adevărat și i s-a arătat lui Simon!”

\VS{35}Ei au povestit cele întâmplate pe drum și cum a fost recunoscut de ei la frângerea pâinii.


\par }{\PP \VS{36}Pe când ziceau ei aceste lucruri, Isus însuși a stat în mijlocul lor și le-a zis: “{\WJ{Pace vouă!”
}}


\par }{\PP \VS{37}Dar ei, îngroziți și înfricoșați, au crezut că au văzut un duh.


\par }{\PP \VS{38}El le-a zis: “{\WJ{De ce vă tulburați? De ce se ridică îndoieli în inimile voastre?
}}

\VS{39}{\WJ{Vedeți mâinile și picioarele Mele, că sunt cu adevărat Eu. Atingeți-mă și vedeți, pentru că un duh nu are carne și oase, așa cum vedeți că am eu.”
}}

\VS{40}După ce a spus acestea, le-a arătat mâinile și picioarele.

\VS{41}Pe când ei încă nu credeau de bucurie și se mirau, le-a zis:
{\WJ{“Aveți ceva de mâncare aici?”
}}


\par }{\PP \VS{42}I-au dat o bucată de pește fript și un fagure de miere.

\VS{43}El le-a luat și a mâncat în fața lor.

\VS{44}El le-a zis: “{\WJ{Iată ce v-am spus când eram încă cu voi: trebuie să se împlinească tot ce este scris în Legea lui Moise, în profeți și în psalmi despre Mine.”
}}


\par }{\PP \VS{45}Atunci le-a deschis mintea, ca să înțeleagă Scripturile.

\VS{46}Și le-a zis:
{\WJ{“Așa este scris și așa a fost necesar ca Hristosul să sufere și să învieze din morți a treia zi,
}}

\VS{47}{\WJ{iar în numele Lui să se propovăduiască pocăința și iertarea păcatelor la toate neamurile, începând de la Ierusalim.
}}

\VS{48}{\WJ{Voi sunteți martori ai acestor lucruri.
}}

\VS{49}{\WJ{Iată, Eu trimit asupra voastră făgăduința Tatălui Meu. Dar așteptați în cetatea Ierusalimului până când veți fi îmbrăcați cu putere de sus.”
}}


\par }{\PP \VS{50}I-a condus până în Betania, și, ridicându-și mâinile, i-a binecuvântat.

\VS{51}În timp ce-i binecuvânta, s-a retras de la ei și a fost dus la cer.

\VS{52}Ei, închinându-se lui, s-au întors la Ierusalim cu mare bucurie,

\VS{53}și stăteau neîncetat în templu, lăudând și binecuvântând pe Dumnezeu. Amin.


\par }